\section{The dimension of the common cuspidal space}\label{sec:count}

The off-diagonal Gram estimate is measured relative to the dimension of a common target. Its level and archimedean types are already prescribed by the periods. We compute the trace of the projector onto precisely these data, so its identity term has the same normalization as the Gram sum. Bounding the unipotent
and boundary terms relative to the identity term gives the
dimension estimate with rank, level and field degree allowed
to vary. This estimate normalizes the selected-row Gram
bound in the next chapter. Definition~\ref{hyp:cuspidal-dimension} specifies the convolution idempotent and the discrete-series parameter at every infinite place. The dimension in \eqref{eq:dimension-count} refers to their simultaneous image.

The following trace calculation draws on companion expositions of
the dimension of the common cuspidal space. \cite{companion-row-count}
\begin{paramdefinition}[Level, signature and weight gap for the trace estimate]\label{hyp:cuspidal-dimension}
Let \(r=p+q\) be the ambient unitary rank. At every real place the signature is \((p,q)\), where \(p\ge q\) and \(q^2\le r\). The CM extension
\(E/E^+\) is unramified at the finite places in this comparison.
Write \(f=[E^+:\Q]\). Both \(r\) and \(f\) may grow.

Let \(g\) be the level
size, \(g_0>4r\) a rational prime split completely in \(E\)
and unramified there, and \(d_E\) the root discriminant of \(E\).
The finite projector \(e\) is a self-adjoint convolution idempotent
on the lattice stabilizer \(K_f\), bi-invariant under its principal
level subgroup \(K_g\). Its support reduces to unipotents at \(g_0\).
For the Harish--Chandra parameter \(\lambda_v\), let \(I_+, I_-\)
be the two signature blocks of the compact Cartan. Require, at every
real place,
\[
 |\lambda_{v,i}|\le b^2\quad(1\le i\le r),\qquad
 |\lambda_{v,i}-\lambda_{v,j}|\ge h\quad(i\in I_+,\ j\in I_-).
\]
The second condition is vacuous when \(q=0\).
Assume
\[
 \begin{gathered}
 1\le g\le b^{C_g},\qquad b^2\ge r^{19},\qquad
 \log d_E\le e^r,\\
 \log f\le\sqrt r,\qquad \log\log b\le rf\log r/16,\\
 h\ge 8\cdot10^{14}(C_g+4)\,r^{19}f^5\log(bd_E).
 \end{gathered}
\]
Here \(C_g\) is fixed and \(r\) is sufficiently large. \end{paramdefinition}
The instance in the ambient rank count is checked in Proposition~\ref{prop:selected-rank-comparison}. In particular the exponent for \(g=g_0g_1\) is twice the common exponent for either prime. The split prime giving the support condition is \(g_0\).

The rank-dependent archimedean gap will control the far orbital sum.

For each infinite place \(v\), let \(\pi_v\) be the
discrete series of the prescribed weight, \(w_v\) a unit vector
of its lowest compact type, and \(d_v\) its formal degree for
the chosen Haar measure. Put \(w=\bigotimes_v w_v\) in the
archimedean tensor product \(\pi=\bigotimes_v\pi_v\). We will
count archimedean intertwiners into the image of the finite
projector. Write \(D\) for their dimension and \(V\) for the
quotient volume. The identity term in the trace is
\[
        D_*=\left(\prod_v d_v\right)e(1)V.
\]

\begin{proposition}[Dimension of the projected cuspidal space]\label{prop:projector-dimension}
For the self-adjoint finite projector and the discrete-series weights
of Definition~\ref{hyp:cuspidal-dimension}, the multiplicity \(D\) is finite. For \(q=0\) one has \(D=D_*\). For sufficiently large \(r\),
uniformly over the fields, levels and admissible weights,
\begin{equation}\label{eq:dimension-count}
    |D-D_*|\le e^{-rf\log(r)/16}D_* .
\end{equation}
The threshold depends on the fixed constants in the hypotheses.
\end{proposition}

The argument first bounds the periodized kernel on the moderate
part of the quotient. On its complement we estimate a finite
orthonormal family before taking the trace. This order will also
establish that the projected space has finite dimension.

\subsection{The finite projector and its normalization}

The value of the finite projector at the identity enters both the trace and its main term. Passing to the finite quotient computes this value without fixing the volume of the maximal compact subgroup to be one.

Let \(H=K_f/K_g\), equipped with counting measure, and let
\(\kappa=\operatorname{vol}(K_g)\) for the finite Haar measure.
Write \(e(x)=a(\bar x)/\kappa\) on \(K_f\), so convolution is the counting-measure convolution of \(a\) on \(H\).
\begin{lemma}[Mass of the finite projector]\label{lem:finite-projector-mass}
The finite convolution projector is of positive type, satisfies
\(|e(x)|\le e(1)\), and has
\[
 e(1)=\frac{\operatorname{rank}(a\mid\C[H])}{\operatorname{vol}(K_f)}.
\]
In particular the identity contribution \(D_*\) is unchanged by a
rescaling of any of the chosen Haar measures.
\end{lemma}

\begin{proof}
For a self-adjoint
idempotent \(a\),
\[
  a(1)=\sum_{x\in H}|a(x)|^2,\qquad |a(x)|\le a(1).
\]
The equality is the convolution identity at 1. The bound at any
\(x\) follows by Cauchy--Schwarz from the convolution identity at
\(x\). In addition \(e\) is of positive type on the finite adelic
group. In finite adelic \(L^2\) its values are coefficients of the right
regular action on \(e\):
\[
 \int\overline{e(y)}e(yx)\,dy=(e*e)(x)=e(x),\qquad |e(x)|\le e(1).
\]
This gives the positive-type kernel used in the automorphic count.

The two finite constructions used by the count can be treated in
this description. A unitary character of an open subgroup gives
its normalized averaging idempotent. It is constant on the level
cosets when the character is trivial on \(K_g\). In the two-prime
construction, the first reduction has a subgroup character
idempotent supported on the required unipotent reductions. The
other reduction uses the product of a subgroup character
idempotent with a central idempotent.

In a unitary irreducible block \(\rho\), a central idempotent
\(z\) acts by a scalar \(c_\rho\) satisfying \(c_\rho^2=c_\rho\).
Thus \(c_\rho\in\{0,1\}\) and \(z^*=z\). For the subgroup
character projector \(a_\chi\) the two identities needed at this
reduction are
\[
 (a_\chi z)^*=z a_\chi=a_\chi z,\qquad
 (a_\chi z)^2=a_\chi^2 z^2=a_\chi z .
\]
The tensor of the two factors inflates by division by the Haar
volume of the common reduction kernel. It is bi-invariant under
each level subgroup inside this kernel. The first factor gives
the stated unipotent support at \(g_0\).

In the delta basis of the regular representation of \(H\) every
diagonal entry of convolution by \(a\) is \(a(1)\). Hence
\[
 \operatorname{rk}(a)=\operatorname{tr}(a)=|H|a(1),\qquad
 e(1)=\frac{\operatorname{rk}(a)}{\operatorname{vol}(K_f)} .
\]
For a subgroup character this is the reciprocal of its supporting
subgroup's volume. The independent \(\kappa^{-1}\) in the pointwise
bound comes from unfolding \(K_g\). In \(D_*\), a change in finite
Haar measure scales \(V\) and \(e(1)\) inversely. At a real place
the same cancellation pairs \(V\) with the formal degree.

\end{proof}

\subsection{Decay of the archimedean projector}\label{sec:count-decay}

The finite projector restricts rational elements. At the real places the spectral gaps must give uniform decay away from the compact subgroup.

At one real place use the coordinate order of the strictly decreasing
Harish--Chandra parameter \(\lambda=(\lambda_1,\ldots,\lambda_r)\).
Let \(J\) be the diagonal signature matrix in this order, and
\(K=\mathrm U(p)\times\mathrm U(q)\) its compact centralizer in
\(\mathrm U(r)\). The gap condition is
\(\lambda_i-\lambda_j\ge h\) for \(i<j\) of opposite signs.
This allows either sign at each index in the weight order.

Write \(S\) for the irreducible lowest \(K\)-type, \(w\) for a unit
highest weight vector of \(S\), and \(P_z\) for the rank-one
orthogonal projector onto a unit vector \(z\). For a Hermitian
matrix \(X\) with \(XJ=-JX\), the compression of \(\pi(\exp X)\) to
\(S\) is
\begin{equation}\label{eq:count-compression}
 (\dim S)\int_K \Phi_\lambda(k^{-1}\exp(X)k)P_{\pi(k)w}\,dk .
\end{equation}
The Haar measure on \(K\) has mass one. Let \(\rho=\rho_C+\rho_N\) be the half-sum of the positive roots,
split into compact and noncompact roots.
If \(m_i(A)\) denotes the
leading principal minor of order \(i\) and \(m_0=1\), define
\[
 \Phi_\lambda(A)=\prod_{i=1}^r \left(\frac{m_i(A)}{m_{i-1}(A)}\right)^{-\mu_i},
 \qquad \mu=\lambda+\rho.
\]
The exponents \(\mu_i\) are
integers by the Harish--Chandra parameter condition. For \(A\) in
\eqref{eq:count-compression}, \(m_r(A)=1\) since
\(\operatorname{tr} X=0\). Thus only the powers
\(-\mu_i+\mu_{i+1}=-(\lambda_i-\lambda_{i+1}+1)\) for \(i<r\) remain.

\begin{proposition}[Lowest-type compression and polar decay]\label{prop:lowest-type-polar-decay}
Let \(\pi\) be the discrete series of
Definition~\ref{hyp:cuspidal-dimension}, with lowest \(K\)-type \(S\),
and let \(X=X^*\) satisfy \(XJ=-JX\). For \(q\ge1\), the lowest-type compression is the flag integral
\eqref{eq:count-compression}. Writing \(P_S\) for the lowest-type projector, its norm bound \eqref{eq:count-polar-decay} is 
\[
 \|P_S\pi(\exp X)|_S\|_{\rm op}\le\exp\!\left[-\frac{h}{8}\min\!\left(s^2,\frac{s}{\sqrt{2q}}\right)\right],\qquad s=\|X\|_{\rm HS}.
\]

If \(q=0\), the relation \(XJ=-JX\) forces \(X=0\) and the
compression is the identity.
\end{proposition}

\begin{proof}
\proofstep{The compact flag integral.}
Let \(H(A)\) denote the expression on the right in
\eqref{eq:count-compression}, with \(A\) in place of \(\exp X\).
It is defined on positive definite matrices and is holomorphic
where all the minors in the integral are uniformly nonzero,
in particular near the identity. Write \(\tau\) for the
holomorphic extension of the action on \(S\) to the complexified
compact group \(K_{\C}\). It satisfies \(\tau(a^*)=\tau(a)^*\),
where the matrix star in this paragraph denotes Hermitian transpose.
This holds on \(K\) by unitarity and extends to \(K_{\C}\) by
holomorphic continuation. The highest weight and the minor exponent satisfy
\[
 \Lambda=\lambda+\rho_N-\rho_C,\qquad
 \mu=\lambda+\rho_N+\rho_C=\Lambda+2\rho_C .
\]
This difference of weights determines the Jacobian in the flag integral.

For \(a\in K_{\C}\) and \(k\in K\), apply QR decomposition in each
sign space using the induced coordinate order. Write \(ak=cb\),
where \(c\in K\) and \(b\) is upper triangular with positive diagonal
\(d_i\). Write \(b_i\) for the first \(i\)-by-\(i\) block of \(b\), and
\(A_i\) for that of \(A\). Triangularity implies
\[
 (b^*Ab)_i=b_i^* A_i b_i,\qquad
 m_i(b^*Ab)=|\det b_i|^2m_i(A)=\prod_{j\le i}d_j^2\,m_i(A).
\] Since the upper unipotent fixes the highest weight vector, the
minor and vector factors transform by
\[
 \begin{aligned}
 \Phi_\lambda(k^{-1}a^*Aa k)&=d^{-2\mu}\Phi_\lambda(c^{-1}Ac),\\
 \tau(k)w&=d^\Lambda\tau(a)^{-1}\tau(c)w,
 \end{aligned}
 \qquad d^\nu=\prod_i d_i^{\nu_i}.
\] The map of compact flags \(kT\mapsto cT\), for the diagonal
torus \(T\) of \(K\), has real Jacobian \(d^{-4\rho_C}\).

In the tangent spaces translated to the standard flag, its complex
differential is \(\operatorname{Ad}(b)\) on the quotient of the
compact complexified Lie algebra by its upper triangular subalgebra.
The quotient has one direction for each negative compact root, and
the diagonal of the induced triangular action is \(d_j/d_i\) for
\(i<j\) of the same sign. Squaring the absolute value of the
determinant gives the stated Jacobian in the invariant flag measure.
The integral descends to this flag space because a torus element
preserves the highest weight line and conjugation by it preserves
all the leading minors.

Inserting the two weights in the integral
gives precisely \(d^{-2\mu+2\Lambda}=d^{-4\rho_C}\). Therefore
\[
       H(a^* A a)=\tau(a)^{-1}H(A)\tau(a^*)^{-1} .
\]
Near identity this identity extends holomorphically in \(A\) and in
the independent variables replacing \(a^*\) and \(a\). Indeed the
Hermitian matrices and the pairs \((a^*,a)\) give real forms in
local complex coordinates for those variables. The local holomorphic identity is therefore
\[
 H(uAv)=\tau(v)^{-1}H(A)\tau(u)^{-1}
       \qquad(u,v\in K_{\mathbb C}\text{ near }I).
\]
It follows that the germ \(g\mapsto H(g^*)\) on the real unitary
group has the left and right \(K\)-covariance of the compression.
Its value at identity is \(I_S\) by compact Schur orthogonality.

\proofstep{The first-order equations.}
The complexified noncompact tangent space \(\mathfrak p_{\C}\) has
basis \(E_{ij}\) with opposite signs. Its tensor product with \(S\)
splits into two orthogonal sums \(V_+\) and \(V_-\) of compact types
with highest weights \(\Lambda+\alpha\) for positive and negative
noncompact roots \(\alpha\), respectively. To see the list of types,
let \(W_+,W_-\) be the two sign spaces. Under the complexified
compact group,
\[
 \mathfrak p_{\C}=(W_+\otimes W_-^*)\oplus(W_-\otimes W_+^*).
\]
The matrix unit from the \(j\)-th to the \(i\)-th coordinate
has weight \(e_i-e_j\), with the signs of \(i,j\) opposite. The one-box Pieri rule and its dual add a single
coordinate in the first factor and subtract a single coordinate in
the other, retaining exactly the dominant results. These are the
indicated noncompact roots.

In the discrete series every compact type
has highest weight \(\Lambda\) plus a sum of positive roots for the
order of \(\lambda\), by Atiyah--Schmid, Theorem 8.5. Thus the action
map \(\mathfrak p_{\C}\otimes S\to\mathcal H_\pi^\infty\) kills
\(V_-\). This gives the first-order equations for the compression
\(g\mapsto P_S\pi(g)|_S\). For \(v_-=\sum E_{ij}\otimes u_{ij}\in V_-\), differentiate
the compression \(F(g)=P_S\pi(g)|_S\) in the right directions:
\[
 \sum D_{ij}F(g)u_{ij}
   =P_S\pi(g)\left(\sum d\pi(E_{ij})u_{ij}\right)=0 .
\]
The vanishing is precisely the absence of \(V_-\) in the
discrete-series action on \(\mathfrak p_{\C}\otimes S\).

For an upper strictly triangular \(Y\), right multiplication
by \(\exp(tY)\) acts on each initial block of columns by an
upper unitriangular matrix. In particular,
\[
 m_l(Z\exp(tY))=m_l(Z)\quad(1\le l\le r),\qquad D_Y\Phi_\lambda=0 .
\] The scalar in \(H(g^*)\) is
\(\overline{\Phi_\lambda(k^{-1}gk)}\). On the real Lie algebra
\(\overline Y=-JY^tJ\), so after complexification and conjugation
by \(k\) its derivative only uses the original entries \(Y_{ij}\)
with \(i\le j\). The conjugate of each vanishing upper-direction
derivative is a vanishing lower-direction derivative.
In the tensor of differentiation directions and input vectors this
conjugation is the compact action by \(k^{-1}\), which preserves
\(V_-\). For \(i<j\), the component indexed by the matrix unit \(E_{ij}\) has zero coefficient along \(w\). Otherwise that tensor would have a weight component
\(\Lambda+e_i-e_j\).

Indeed every weight of a constituent with highest weight
\(\Lambda-\beta\), for a positive noncompact root \(\beta\), has the form
\[
 \Lambda-\beta-\sum_{\alpha\in\Delta_C^+} n_\alpha\alpha,
       \qquad n_\alpha\in\mathbb N ,
\]
where \(\Delta_C^+\) is the set of positive compact roots. This
omits \(\Lambda+e_i-e_j\) for \(i<j\), since a nonempty sum of
positive roots cannot vanish. 

The diagonal matrix units do not belong to \(\mathfrak p_{\C}\). After conjugation by \(k\), write the input tensor as \(\sum_{i,j} E_{ij}\otimes u_{ij}(k)\), summing over the noncompact root spaces. The differentiated integrand
is a sum of scalar derivatives times
\(\langle w,u_{ij}(k)\rangle\pi(k)w\).
If \(a_{ij}(k,g)\) denotes the scalar derivative in the direction \(E_{ij}\), the
two tests give
\[
 a_{ij}(k,g)=0\quad(i>j),\qquad
 \langle w,u_{ij}(k)\rangle=0\quad(i<j).
\]
There is no diagonal term, so each summand
\(a_{ij}(k,g)\langle w,u_{ij}(k)\rangle\pi(k)w\) vanishes. Thus the first-order equation holds before integration.

\proofstep{Determination from the value at identity.}
Let \(F\) be the difference of the two germs. It is zero at identity.
Suppose derivatives of all orders less than \(m\) vanish there.
At order \(m\) any compact differentiation can be commuted to the
right, with commutators of lower order. Right covariance then
reduces it to order \(m-1\). Interchanging two noncompact directions produces a commutator of
lower order, which vanishes by induction. Thus the remaining derivative belongs to
\[
 \operatorname{Hom}_K(\operatorname{Sym}^m\mathfrak p_{\mathbb C}
                         \otimes S,S).
\] Differentiating the
first-order equations shows the map kills tensors with \(V_-\)
in the product of any one noncompact tensor factor with \(S\). It factors through
the orthogonal complement of the sum of these subspaces. In this
complement each noncompact factor paired with \(S\) lies in \(V_+\).

Let \(c_S\) be the compact Casimir eigenvalue on \(S\), in the trace
form. On the tensor product the compact Casimir is
\(c_S+\Omega_{\mathfrak p^{\otimes m}}+2\sum_i M_i\), where
\(M_i\) is the mixed term between the \(i\)-th tensor factor and \(S\).
Compact unitarity gives
\[
 \langle u,\Omega_{\mathfrak p^{\otimes m}}u\rangle\ge0 .
\] The Casimir on a single \(\mathfrak p_{\C}\)
factor equals \(p+q=r\), since the standard and dual standard
representations of \(\mathrm U(d)\) each have trace-form Casimir
\(d\). Thus on a positive type \(\Lambda+e_a-e_b\), \(2M_i\)
has eigenvalue
\[
  2(\lambda_a-\lambda_b+\rho_N(a)-\rho_N(b))+2-r
                 \ \ge\ 2(h-2r).
\]
Here \(|\rho_N(a)|\le r/2\).

Since \(h>2r\) in the noncompact
case, the quadratic form of the tensor Casimir on the complement
is strictly larger than \(c_S\) times the squared norm. The
complement is \(K\)-invariant, so it has no summand isomorphic to
\(S\). This forces the derivative map to vanish and completes the
induction.

Along \(g=\exp(tX)\) with \(X\) Hermitian and anticommuting with
\(J\), the integral is real analytic for all real \(t\). All the
leading minors are positive there and compactness gives a common
complex neighborhood of each \(t\) for the integration. The
compression is real analytic on this curve as well. In fact the
infinitesimal norm estimate \eqref{eq:count-step} below at one
place grows at most linearly in the word length. Iteration from the lowest type gives a constant \(C_X\) such that
\[
 \left\|\frac{d^m}{dt^m}\pi(\exp(tX))u\right\|
       \le m! C_X^m\|u\|\qquad (m\ge0,\ t\in\mathbb R).
\]
Unitarity makes the bound uniform in \(t\), so Taylor's estimate
gives analyticity. The vanishing
at identity of all derivatives therefore proves
\eqref{eq:count-compression} for every such \(X\). The quantitative threshold \(h\ge120rq^{3/2}\) implies
\(h>2r\) at every noncompact place.

\proofstep{Quantitative decay.}

The local expression for the compressed coefficient is now explicit. Its useful decay must be expressed in terms of the opposite-sign spectral gap, which remains large throughout the count.


Put \(Y=k^{-1}Xk\) and \(s=\|Y\|_{\mathrm{HS}}=\|X\|_{\mathrm{HS}}\).
On the \(i\)-th exterior power let \(L_i(Y)\) act by applying
\(Y\) to each wedge factor and summing. With
\(e_{[i]}=e_1\wedge\cdots\wedge e_i\), the minor has the spectral representation
\[
 m_i(\exp Y)=\langle e_{[i]},\exp(L_i(Y))e_{[i]}\rangle.
\]

The spectral measure of \(L_i(Y)\) at this unit vector is a
probability distribution \(\xi_i\). Its mean is zero because all
diagonal entries of \(Y\) vanish. Its variance is
\[
   V_i=\|L_i(Y)e_{[i]}\|^2=\sum_{a\le i<b}|Y_{ba}|^2.
\]

Indeed, for each term off the diagonal, the replacement of the
\(a\)-th factor by \(e_b\) either repeats a factor or gives one of
the orthonormal wedges with exactly one index outside \([i]\).
The eigenvalues of \(Y\) are pairs \(\pm t_j\) with \(t_j\ge0\),
and zeros, by the singular-value decomposition of the block between
the two sign spaces. There are at most \(q\) pairs and
\(2\sum_j t_j^2=s^2\). Each eigenvalue of \(L_i(Y)\) is a sum drawn without replacement
from this multiset. In particular
\[
              \|L_i(Y)\|_{\rm op}\le \sum_j t_j\le\sqrt{2q}\,s .
\]

For \(s>0\), set \(t=\min(1,1/(\sqrt{2q}\,s))\). If \(|z|\le1\),
Taylor's integral remainder and \(e^{uz}\ge1-u\) for \(0\le u\le1\)
give \(e^z\ge1+z+z^2/3\). Apply this to \(z=t\xi_i\). Since
\(t^2V_i\le1\) and \(\log(1+a/3)\ge a/4\) for \(0\le a\le1\),
the log moment satisfies \(\log\mathbb E e^{t\xi_i}\ge t^2 V_i/4\).
Its value at zero is zero, so convexity on \([0,1]\) gives
\[
 \log m_i(\exp Y)=\log\mathbb E e^{\xi_i}
       \ge t^{-1}\log\mathbb E e^{t\xi_i}\ge tV_i/4 .
\] The weighted sum of
the variances satisfies
\[
 \sum_{i<r}(\lambda_i-\lambda_{i+1}+1)V_i
 = \sum_{a<b}(\lambda_a-\lambda_b+b-a)|Y_{ba}|^2
 \ge h s^2/2 .
\]
All nonzero entries in the last sum cross the two signs. Thus the
scalar multiplier in \eqref{eq:count-compression} is bounded above by
\begin{equation}\label{eq:count-polar-decay}
      \exp\!\left[-\frac h8\min\!\left(s^2,\frac{s}{\sqrt{2q}}\right)\right].
\end{equation}

It is nonnegative, and the same bound is immediate for \(s=0\).
By Schur orthogonality,
\((\dim S)\int_K P_{\pi(k)w}\,dk=I_S\). Consequently the multiplier
bound is also the operator norm bound for the compression, with no
factor depending on \(\dim S\). The polar decomposition \(g=k_0\exp X\)
preserves this bound since \(\pi(k_0)\) acts unitarily on \(S\).

\end{proof}

\subsection{The initial bound in the deep part}\label{sec:count-sobolev}

Near a very short lattice, many translates can contribute at a single point. A pointwise bound for any finite orthonormal family is needed before one knows that the whole projected space has finite dimension.

\begin{proposition}[Initial point evaluation in a fixed order]\label{prop:initial-point-evaluation}
Fix the archimedean representation, Haar gauges and finite level of
Definition~\ref{hyp:cuspidal-dimension}, with \(r\ge8\).
Set \(N=r^2 f\), \(B=b_{P_0+N}\) as in \eqref{eq:count-step}, and
let \(C(y)\) be the multiplicity factor in \eqref{eq:gram-chart-multiplicity}.
For a finite family of \(D\) isometric intertwiners and a vector \(u\)
of order at most \(P_0\) at each place, one has
\[
 \sum_{j=1}^D |F_j(u,y)|^2\le (4B)^N C(y)D\|u\|^2 .
\]
For \(P_0=8qN\) this gives the starting bound, uniform on the
neighborhood of a deep point with least size \(\tau(x)\), used
in the descent.
\end{proposition}

\begin{proof}
Use the same isometric intertwiners that give the diagonal kernel.
Write \(F_j(u,y)\) for the continuous automorphic representative at
\(y\) of a vector \(u\) in the \(j\)-th copy of the representation.
It is invariant on the right by the finite level subgroup \(K_g\),
and \(\|F_j(u,\cdot)\|_2=\|u\|\). At each real place take the
filtration by the span of words of at most \(P\) infinitesimal
operators on the lowest \(K\)-type. Use the tensor product of these
spaces, with the same bound \(P\) at each place, as the space \(V_P\).
For a unit Hilbert--Schmidt Lie algebra element at one place,
differentiation sends \(V_P\) into \(V_{P+1}\) with norm at most
\begin{equation}\label{eq:count-step}
 b_P=\sqrt r\,b^2+\tfrac32 r^{3/2}+3P .
\end{equation}

To see this, let \(\rho_C,\rho_N\) be the half-sums of positive
compact and noncompact roots, and \(\rho=\rho_C+\rho_N\). The
lowest \(K\)-type has highest weight \(\Lambda=\lambda+\rho_N-\rho_C\).
For the derived group this is the discrete-series lowest-type formula
of Atiyah and Schmid, Theorem 8.5. The central character supplies
the common central coordinate when passing from
\(\mathrm{SU}(p,q)\times \mathrm U(1)\) to its finite central quotient
\(\mathrm U(p,q)\).
The infinitesimal character is \(\lambda\) in the normalization in
which the quadratic Casimir for the trace form acts by
\(\|\lambda\|^2-\|\rho\|^2\). \cite{atiyah-schmid}
The lowest isotypic component consists of smooth vectors.

Write \(P_S\) for the compact isotypic projection. Compact
averaging preserves smoothness, and continuity gives
\[
 P_S(\mathcal H_\pi^\infty)\subseteq S\cap\mathcal H_\pi^\infty,
 \qquad \overline{P_S(\mathcal H_\pi^\infty)}=P_S(\mathcal H_\pi)=S .
\]
Since \(S\) is finite dimensional, the projected dense linear space
equals \(S\).
For the word space \(V_P\), the surjection
\[
 \bigoplus_{j=0}^P\mathfrak g_{\C}^{\otimes j}\otimes S
      \longrightarrow V_P,\qquad
 X_1\otimes\cdots\otimes X_j\otimes u
       \longmapsto d\pi(X_1)\cdots d\pi(X_j)u
\]
is \(K\)-equivariant.
Thus each occurring compact highest weight \(\nu\) is a weight of
the first factor plus at most \(P\) roots of length \(\sqrt2\).
The weights of the first factor lie in the convex hull of the compact
Weyl orbit of \(\Lambda\). For any compact Weyl permutation \(w\), dominance gives
\[
 \|\Lambda+\rho_C\|^2-\|w\Lambda+\rho_C\|^2
      =2\langle\Lambda-w\Lambda,\rho_C\rangle\ge0 .
\]
Convexity extends the norm bound to every weight in the convex hull. Adding the contributions from the lowest type and the \(P\) adjoint
factors yields
\[
         \|\nu+\rho_C\|\le\|\lambda\|+\|\rho_N\|+\sqrt2 P .
\]

Let \((E_i)\) be a Hilbert--Schmidt orthonormal basis
adapted to the compact and noncompact directions. In the trace
form the compact basis contributes with a minus sign and the
noncompact basis with a plus sign. Thus in the enveloping algebra
\[
 \Omega_K=-\sum_{i\ {\rm compact}}E_i^2,\qquad
 \Omega_G=\Omega_K+\sum_{i\ {\rm noncompact}}E_i^2.
\]
Since \(d\pi(E_i)\) is skew-adjoint on smooth vectors, it follows
that \(\sum_i\|d\pi(E_i)u\|^2
=\langle u,d\pi(2\Omega_K-\Omega_G)u\rangle\). The operator \(\Omega_K\) uses the trace form
restricted to the compact subalgebra and acts on a compact type
\(\nu\) by \(\|\nu+\rho_C\|^2-\|\rho_C\|^2\).

With
\(A=\|\lambda\|\) and \(\delta=\|\rho_N\|+\sqrt2 P\), the sum is
therefore at most \((2(A+\delta)^2-A^2+\|\rho\|^2)\|u\|^2\).
In fact
\[
 2(A+\delta)^2-A^2+\|\rho\|^2
        \le (A+2\delta+\|\rho\|)^2,\qquad
 A+2\delta+\|\rho\|
        \le \sqrt r\,b^2+\tfrac32r^{3/2}+2\sqrt2 P\le b_P .
\]
The second comparison uses the bound \(r/2\) on each coordinate
of \(\rho,\rho_N\) and \(b^2\) on \(|\lambda_i|\). For an arbitrary unit direction use
Cauchy--Schwarz in the orthonormal basis, and tensor the resulting
bound at a single place with the identity at the others.

\proofstep{An injective product chart and its Jacobian.}

Write \(N=r^2 f\) and \(B=b_{P_0+N}\) for a fixed starting order
\(P_0\). Take \(a=1/B\). By Definition~\ref{hyp:cuspidal-dimension}, \(b^2\ge r^{19}\);
for \(r\ge8\) the chart radius therefore satisfies
\[
   a=B^{-1}\le\min\{(8r^5)^{-1},(12r^3)^{-1}\}.
\]
At each real place use an orthonormal Lie algebra basis \(E_i\)
and the coordinates
\(t\mapsto\Theta(t)=\prod_{i=1}^{r^2}\exp(t_i E_i)\) on
\([0,a]^{r^2}\). For a product \(U=\prod_j\exp(s_j E_{i_j})\), put \(L=\sum_j|s_j|\); negative parameters include inverses. Expanding the product and using \(\|E_i\|_{\rm op}\le1\) gives 
\[
 \|U\|_{\rm op}\le\prod_j e^{|s_j|}=e^L,\qquad \|U-I\|_{\rm op}\le\prod_j e^{|s_j|}-1=e^L-1.
\]

For the chart differential, its \(i\)-th column is a product of a
prefix, \(E_i\), and a suffix. The ambient Hilbert--Schmidt distance
from \(E_i\) is at most \(e^{r^2 a}-1\le 2r^2a\). Let \(S_i=\prod_{j>i}\exp(t_jE_j)\). Left translation gives
\[
 \Theta(t)^{-1}\partial_i\Theta(t)=S_i^{-1}E_i S_i,\qquad
 \|S_i^{-1}E_iS_i-E_i\|_{\rm HS}\le e^{2r^2a}-1\le3r^2a .
\]
Indeed subtract \(E_iS_i\) and use the bounds on \(S_i^{\pm1}\)
together with \(\|E_i\|_{\rm HS}=1\).
The operator norm of a map with \(r^2\) such error columns is at
most \(r\) times the maximum column norm. Set \(I_E(t)=\sum_i t_iE_i\), the linear map into the
ambient matrix space. For \(t,t'\) in the cube, integration of
the ambient differential along the segment gives
\[
 \|\Theta(t)-\Theta(t')-I_E(t-t')\|_{\rm HS}
                         \le 2r^3a\|t-t'\| .
\]

Since \(\|I_E(t-t')\|_{\rm HS}=\|t-t'\|\)
and \(2r^3a<1\), the chart is injective. For the
left-translated differential each singular value is at least
\(1-3r^3a\), so the Jacobian with respect to the orthonormal basis satisfies
\[
 |\det d\Theta|\ge(1-3r^3a)^{r^2}\ge e^{-6r^5a}\ge e^{-1}.
\]
Here \(3r^3a\le1/4\) and \(6r^5a\le1\).
Let \(c_v\) be the reciprocal of the density of Haar volume at the
identity relative to Hilbert--Schmidt Lebesgue measure. Write \(dt_v\) for coordinate Lebesgue measure and \(dg_v\) for
Haar measure. On the image the Jacobian estimate gives
\[
 dt_v\le e c_v\,dg_v,\qquad
 \prod_v dt_v\le e^f\Big(\prod_v c_v\Big)\prod_v dg_v .
\] Two elements of the image
have quotient of operator norm \(\le e^{2r^2 a}\le 2\) at every place.

\proofstep{Unfolding the \(L^2\) integral.}

Write \(\tau_1(y)=\min(1/2,\tau(y))\) for the truncated least
lattice size at \(y\). In unfolding the chart about \(y\) with the
finite coset \(y_fK_g\), two points mapping to the same automorphic
point differ by a rational group element preserving the lattice at
\(y\). Relative to a chosen representative of the fibre, its
archimedean matrix is a quotient of two chart points. Thus
\[
 \|\gamma_v\|_{\rm HS}\le\sqrt r\|\gamma_v\|_{\rm op}
     \le\sqrt r e^{2r^2a}\le2\sqrt r
\]
in the lattice frame at each place.
The lattice separation and product packing in the far orbital count
below bound the number in the fibre by
\(M(y)=(1+4\sqrt r/\tau_1(y)^2)^{2N}\).
Indeed the difference of two distinct elements is a nonzero integral
endomorphism, so the geometric mean of its Hilbert--Schmidt norms
is at least \(\tau_1(y)^2\).

Product balls about the elements with
radius \(b_0=\tau_1(y)^2/2\) at each place thus have disjoint interiors.
In ambient real dimension \(2N\) their volume ratio is bounded by
\[
 \left(\frac{2\sqrt r+b_0}{b_0}\right)^{2N}
     =\left(1+\frac{4\sqrt r}{\tau_1(y)^2}\right)^{2N}.
\] The finite coset contributes its full
volume \(\operatorname{vol}(K_g)\) in the unfolded integral, since
the squared absolute value is right \(K_g\)-invariant. The integral
of \(|F_j(u,\cdot)|^2\) on the product coordinate cube is therefore
at most \(C(y)\|u\|^2\), where
\begin{equation}\label{eq:gram-chart-multiplicity}
 C(y)=e^f\Bigl(\prod_v c_v\Bigr)\frac{M(y)}{\operatorname{vol}(K_g)} .
\end{equation}
This estimate uses the isometry and applies to every
\(u\in V_{P_0+N}\).

\proofstep{Iteration of the point evaluation inequality.}

The evaluation bound has so far been obtained at a fixed derivative order. Iteration allows the depth of the thin region to determine how much decay is needed.


For a differentiable function \(\phi\) on \([0,a]\), the fundamental
theorem of calculus and Cauchy--Schwarz give
\[
 |\phi(0)|^2\le\frac{2}{a}\int_0^a|\phi(t)|^2\,dt
                     +2a\int_0^a|\phi'(t)|^2\,dt .
\]
Start at the origin and append the coordinates in order. At stage
\(i\) fix the prefix \(x=y\exp(t_1 E_1)\cdots
\exp(t_{i-1}E_{i-1})\) and one of the fixed vectors \(w\) from
the preceding stages.

Apply the one-variable inequality to
\(t\mapsto F_j(w,x\exp(tE_i))\). Right differentiation inserts the next word operator:
\[
 \frac d{dt}F_j(w,x\exp(tE_i))=F_j(d\pi(E_i)w,x\exp(tE_i)).
\] The two resulting
cube integrands thus use the fixed vectors \(w\) and
\(d\pi(E_i)w\), with weights \(2/a\) and \(2a\),
respectively. Applying the next inequality at each point of this
prefix leaves these vectors fixed as the next factor is appended.
After \(N\) stages apply the cube bound to each resulting vector.
The sum of their squared norms with these weights is at most
\((2/a+2aB^2)^N\|u\|^2\), so for \(u\in V_{P_0}\)
\[
 \sum_{j=1}^D |F_j(u,y)|^2\le(2/a+2a B^2)^N C(y)D\|u\|^2
                          =(4B)^N C(y)D\|u\|^2 .
\]
Each derivative raises the order by at most one, so \(P_0+N\)
suffices for the whole iteration. For the subsequent descent one may enlarge \((4B)^N\) to \((6B)^N\).

In the unipotent neighborhood used below the finite lattice is
fixed. The archimedean change of frame at each place is
\(\exp(A_v)\exp(Z_v)\), with
\(\|A_v\|_{\mathrm{HS}}+\|Z_v\|_{\mathrm{HS}}\le\rho\le\log 2\).
The matrix and its inverse both have operator norm at most two.
Consequently each covolume in rank \(d\) loses at most a factor
\(2^d\) at one place in the complex Euclidean determinant norm.
For a fixed subquotient of rank \(d>0\) the inverse frame has
\(\|\bigwedge^d g_v^{-1}\|\le2^d\). The induced quotient norm obeys
the same bound by taking the infimum over lifts. Thus
\[
  C_y/C_x\ge\prod_{v\mid\infty}2^{-2d}=2^{-2fd},\qquad
  (C_y/C_x)^{1/(2fd)}\ge\tfrac12
\]
for its normalized covolumes.

Taking the least and truncating at \(1/2\) shows
\(\tau_1(y)\ge\tau(x)/2\) when \(\tau(x)<1/2\). There
\(1+4\sqrt r/\tau_1(y)^2\le17\sqrt r/\tau(x)^2\). The factor for
that neighborhood follows with \(P_0=8qN\). These bounds use the
given lattice and gauge and hold uniformly in \(r,f,b,P_0\) under
the stated size conditions.
\end{proof}

\subsection{The cuspidal projection}\label{sec:count-projection}

A sum of translated matrix coefficients has the desired trace interpretation only after its action on cuspidal forms is identified. Reproduction on the chosen compact type is the local step in this identification.

\begin{auxlemma}[A compactly supported reproducing function]
\label{lem:compact-type-reproducer}
Let \(G\) be a real Lie group, \(K\) a compact subgroup, and \(\pi\)
a strongly continuous unitary representation of \(G\). Suppose the
irreducible \(K\)-type \(S\) occurs with multiplicity one.
For any neighborhood of the identity there is a smooth function
\(\alpha\) supported compactly in that neighborhood with
\(\pi(\alpha)|_S=I_S\). For \(X\in\mathfrak g\) and \(w\in S\),
\[
 d\pi(X)w=\pi(\alpha_X)w,\qquad
 \alpha_X(g)=\left.\frac{d}{dt}\alpha(\exp(-tX)g)\right|_{t=0}.
\]
\end{auxlemma}
\begin{proof}
Compactness of \(K\) gives a smaller neighborhood whose
\(K\)-conjugates all lie in the prescribed one. Average a smooth,
nonnegative approximate identity supported there over conjugation
by \(K\), using probability Haar measure on \(K\). Conjugation by
\(K\) preserves left Haar measure on \(G\), since its modular
character is trivial on \(K\). The averaged functions \(a_j\) thus
have integral one and \(\pi(a_j)\) commutes with \(K\). Schur's lemma
and multiplicity one give
\[
 \pi(a_j)|_S=c_j I_S,\qquad
 |c_j-1|=\|\pi(a_j)w-w\|\longrightarrow0\quad(\|w\|=1,\ w\in S).
\]
Choose \(c_j\ne0\) and set \(\alpha=a_j/c_j\). In
\(\pi(\exp tX)\pi(\alpha)w\), left invariance moves \(\exp(tX)\)
from the representation to the argument of \(\alpha\).
Differentiation under this compactly supported integral gives the
formula.
\end{proof}

\begin{proposition}[Cuspidal projection at the given level]
For the finite idempotent of Definition~\ref{hyp:cuspidal-dimension}
and the rank-one projectors onto the unit lowest vectors at infinity, the resulting global kernel projects
onto the corresponding cuspidal copies with their isometric
Petersson realizations. If \(f\) is one of the smooth right-level-invariant vectors used in the kernel, its constant term along any proper parabolic with unipotent radical \(U\) vanishes over the totally real base \(F=E^+\): 
\[
 \int_{U(F)\backslash U(\mathbb A_F)} f(ug)\,du=0\qquad\text{for every }g.
\]
 The same holds for right translates and differentiated vectors.
\end{proposition}

\begin{proof}
\proofstep{The integrable real coefficient.}

First consider a noncompact real group \(G_v=U(p,q)\).
Write \(D_v(g)=\log\|g\|_{\rm op}\), which is the largest
polar coordinate. This length is symmetric under inversion and
subadditive. Proposition~\ref{prop:lowest-type-polar-decay} gives
\[
   |\langle w_v,\pi_v(g)w_v\rangle|
      \le C e^{-cD_v(g)},\qquad c=\frac{h}{8\sqrt{2q}}.
\]
The bound follows directly for \(D_v(g)\ge1\) and, after increasing
\(C\), for \(D_v(g)\le1\). The gap assumption gives \(c>100r^2\)
at all the places in question. The constants in this paragraph
serve only to prove convergence and vanishing.

Let \(N_v\) be the unipotent radical of the stabilizer of a nonzero
isotropic subspace of dimension \(j\). With respect to this
subspace, a dual isotropic subspace and their common orthogonal
complement, \(N_v\) is block upper triangular with three blocks
and identity diagonal. The logarithm truncates after the quadratic term:
\[
             (n-I)^3=0,\qquad \log n=(n-I)-\tfrac12(n-I)^2 .
\] In logarithmic coordinates Haar
measure is Lebesgue measure by the two-step group law. The
set \(D_v(n)\le R\) has volume \(O(e^{2r^2R})\), since its
logarithms have norm \(O(e^{2R})\) in a real space of dimension at
most \(r^2\).

On the full group, the polar integration Jacobian
is a product of hyperbolic sines of restricted roots, with
multiplicities of total size at most \(r^2\). Indeed the compact
directions transverse to the split diagonal give a factor
\(\sinh\alpha(H)\) per real root direction. On the ball of radius \(R\), each positive root is at most \(2R\).
Writing \(G_R,N_R\) for the two balls gives
\[
  \operatorname{vol}G_R\ll (1+R)^q e^{2r^2R},\qquad
  \operatorname{vol}N_R\ll e^{2r^2R}.
\]
Hence \(e^{-(c/3)D_v}\) is integrable on both groups, since
\(c/3>2r^2\).

Suppress \(v\), and take inner products linear in the second
variable. Let \(\mathcal V\) be the span of the translates of
\(w\). A coefficient with both vectors in \(\mathcal V\) decays
at rate \(c\) by subadditivity. We will also use this rate for
\(\langle v,\pi(g)d\pi(X)w\rangle\), for \(v\in\mathcal V\).
Apply Lemma~\ref{lem:compact-type-reproducer} to the lowest type.
For \(v=\pi(a)w\) and \(g'\in\operatorname{supp}\alpha_X\),
subadditivity gives \(D_v(a^{-1}gg')\ge D_v(g)-D_v(a)-D_v(g')\).
Consequently
\[
 |\langle v,\pi(g)d\pi(X)w\rangle|
 \le C e^{-cD_v(g)} e^{cD_v(a)}
          \int_G |\alpha_X(g')|e^{cD_v(g')}\,dg' .
\]
Finite linear combinations give the bound on \(\mathcal V\).
The same estimate on compact sets of translates permits
differentiation of the averaged coefficient.

Give \(N\backslash G\) its right invariant quotient measure, with
the Weil integration formula for the chosen Haar measures. For
\(v\in\mathcal V\) and \(u=w\) or \(u=d\pi(X)w\) set
\[
  F_{v,u}(Ng)=\int_N\langle v,\pi(ng)u\rangle\,dn,\qquad
  I(v,v')=\int_N\langle v,\pi(n)v'\rangle\,dn .
\]
Write \(\psi(g)=\langle v,\pi(g)u\rangle\). For the integrals
below absolute convergence follows by bounding
\(|\psi(ng)\psi(g)|\) by a constant times
\(e^{-cD_v(n)/2}e^{-cD_v(g)/2}\), using subadditivity.
Unfold one of the two averages in the squared norm and apply
formal-degree orthogonality to the integral on \(G\). This gives
\[
   \|F_{v,u}\|_2^2
   =\int_N\int_G\psi(ng)\overline{\psi(g)}\,dg\,dn
   =\frac{\|u\|^2}{d} I(v,v).
\]
Here inversion on \(N\) leaves Haar measure fixed.
Taking \(u=w\) proves positivity of \(I\). Polarization gives
the corresponding mixed norm identity.

Let \(a(t)=\exp(tH)\) act by \(e^t,1,e^{-t}\) on the
isotropic subspace, the orthogonal middle block and the dual
subspace, respectively. It is central in the block Levi \(M\).
Conjugation by \(a(t)\) multiplies Haar volumes on \(N\) by
\(\delta(t)=e^{\kappa_N t}\), with
\(|\kappa_N|\le2r^2\) from the weights 1 and 2 on the
unipotent blocks. The operators
\[
   U(t)F(Ng)=\delta(t)^{-1/2}F(Na(t)g)
\]
are unitary on the quotient by the Weil integration formula.
Changing variables \(n'=a(-t)na(t)\) in the defining integral gives
\[
  F_{v,w}(Na(t)g)=\delta(t)F_{\pi(a(-t))v,w}(Ng).
\]
Thus \(U(t)\) preserves the closed span of the \(F_{v,w}\).

The generator is bounded on that closed span. In fact
\(G=NMK\), since \(K\) acts transitively on the isotropic
\(j\)-planes. This transitivity follows by sending the two
orthonormal projection frames of such a plane into matching
frames in the positive and negative spaces. At \(g=nmk\)
the path \(Na(t)g\) equals \(Ng\exp(t\operatorname{Ad}(k^{-1})H)\).
At \(t=0\) differentiate \(F_{v,w}\) along this path. For an
orthonormal real Lie algebra basis \((X_i)_i\) in the
Hilbert--Schmidt metric, the sum of squares of the coefficients
of \(\operatorname{Ad}(k^{-1})H\) is \(\|H\|_{\rm HS}^2\).
Write \(\operatorname{Ad}(k^{-1})H=\sum_i b_i(k)X_i\).
At \(t=0\) the derivative of \(U(t)F_{v,w}\) is
\[
 -\tfrac{\kappa_N}{2}F_{v,w}(Ng)
         +\sum_i b_i(k)F_{v,d\pi(X_i)w}(Ng).
\]
Pointwise Cauchy--Schwarz gives
\[
 \left|\sum_i b_i(k)F_{v,d\pi(X_i)w}(Ng)\right|^2
        \le \|H\|_{\rm HS}^2\sum_i |F_{v,d\pi(X_i)w}(Ng)|^2 .
\]

The norm identity divides each integrated square by
\(\|F_{v,w}\|_2^2\) with ratio
\(\|d\pi(X_i)w\|^2/\|w\|^2\) when that denominator is nonzero.
If it is zero, both squares vanish. The derivative is thus bounded
by \(C_1\|F_{v,w}\|_2\), with \(C_1\) independent of \(v\).
The group law and the invariant span give the same bound at other
\(t\). Integrate in \(t\) to obtain
\(\|(U(t)-1)F\|_2\le C_1|t|\|F\|_2\) and pass to the closed span.
Strong continuity and the spectral theorem for this unitary
group now give a bounded self-adjoint generator. For example,
if \(\nu_F\) is the spectral measure of \(F\) in that closed span,
Fatou's lemma gives
\[
 \int \lambda^2\,d\nu_F(\lambda)
 \le\liminf_{t\to0}\frac{\|(U(t)-1)F\|_2^2}{t^2}
 \le C_1^2\|F\|_2^2.
\]

On the other hand, \(D_v(na(t))\ge|t|\) by the eigenvalues
of its triangular blocks, and
\(D_v(n)\le D_v(na(t))+|t|\). Hence
\(3D_v(na(t))\ge|t|+D_v(n)\).
The mixed norm identity gives
\[
   |\langle F_{v,w},U(t)F_{v,w}\rangle|
      \le C_v e^{-(c/3-r^2)|t|}.
\]
Indeed the bound for \(I(v,\pi(a(-t))v)\) follows by the
coefficient estimate on \(\mathcal V\) at \(na(-t)\), with
the two fixed vectors both equal to \(v\) there. The spectral
measure of \(F_{v,w}\) for the bounded generator has compact
support.

For its scalar spectral measure \(\nu\) use
\(\widehat\nu(t)=\int e^{it\lambda}\,d\nu(\lambda)\).
The exponential bound of rate \(c/3-r^2>0\) gives the
holomorphic function
\[
 H(z)=\frac1{2\pi}\int_{\mathbb R}
                e^{-itz}\widehat\nu(t)\,dt
       \qquad (|\operatorname{Im}z|<c/3-r^2).
\]
Fourier inversion on smooth compactly supported tests shows
that \(H(\lambda)\,d\lambda=d\nu(\lambda)\) on the real line.
Thus \(H\) vanishes on the intervals outside the compact
spectral support. The identity theorem makes \(H=0\).
Hence \(\|F_{v,w}\|_2^2=\nu(\mathbb R)=0\).
Continuity in \(g\) and \(v=\pi(x^{-1})w\) give
\[
     \int_N\langle w,\pi(xny)w\rangle\,dn=0
               \qquad (x,y\in G).
\]
This proves the local constant-term identity for all translates
of the coefficient being used.

Return to the adelic quotient \(Y=G(E^+)\backslash
G(\mathbb A_{E^+})\). Use the real test function
\[
 \phi(g)=d\,\overline{\langle w,\pi(g)w\rangle},
 \qquad g\in G_\infty,
\]
with
\(d=\prod_v d_v\) and the tensor vector \(w\) fixed at the
beginning of the section. The tensor product is irreducible
and has Schur orthogonality with formal degree \(d\).
For finite linear combinations \(v,u\) of translates of \(w\)
write \(k_{v,u}(g)=d\,\overline{\langle v,\pi(g)u\rangle}\).
The coefficient decay gives \(L^1\) for these functions, and
orthogonality gives
\[
  k_{v,u}*k_{v',u'}=\langle u,v'\rangle k_{v,u'},
       \qquad k_{v,u}^*=k_{u,v}.
\]
Therefore \(R(\phi)\) is an orthogonal projection in the
archimedean right representation on \(L^2(Y)\).

If \(F\) is
fixed by it, \(v\mapsto R(k_{v,w})F\) extends with norm
\(\|F\|\) to an intertwiner from \(\pi\), by these identities
and the left-translation identity for \(k_{v,w}\). It sends
\(w\) to \(F\). Conversely any such intertwiner \(T\) is determined by \(T(w)\).
Since \(\pi(\phi)\) projects onto \(\mathbb C w\),
\[
                     R(\phi)T(w)=T(\pi(\phi)w)=T(w).
\]
As \(R(e)\) commutes with the archimedean representation,
the vectors fixed by \(P=R(\phi\otimes e)\) are precisely the
values at \(w\) of the intertwiners with image in \(R(e)L^2(Y)\).

The integral kernel of \(P\) is the periodization
\[
   K(x,y)=\sum_{\gamma\in G(E^+)}
                     (\phi\otimes e)(x^{-1}\gamma y).
\]
To justify this identity on compactly supported test functions,
the sum is uniformly absolutely convergent on compact sets in
both variables. The finite support restricts the entries of
\(\gamma\) to a fixed fractional ideal for \(x,y\) in such
sets. If the polar lengths at the real places are bounded by
\(R_v\), the entry sizes are \(O(e^{R_v})\) at those places.
Packing cells of the fixed entry lattice in the enlarged
boxes gives \(O(\prod_v e^{2r^2R_v})\) eligible matrices.

Choose integral shell endpoints \(R_v\ge0\) at the
noncompact places. There the matrix coefficient is at most
a constant times \(e^{-cR_v}\). Thus on each compact set in
the two quotient coordinates, the sum is bounded by
\[
 C\sum_{(R_v)}\prod_{v\ {\rm noncompact}}e^{-(c-2r^2)R_v},
\]
which converges since \(c>2r^2\). At compact
places the entry boxes are bounded. With convergence established,
the kernel formula follows by unfolding \(Y\) in the convolution.

For a proper rational maximal parabolic let \(N\) be its
unipotent radical. The unipotent quotient
\([N]=N(E^+)\backslash N(\mathbb A_{E^+})\) is compact,
by a central series with successive additive quotients and
compactness of the additive adelic quotient. Average \(K(nx,y)\)
over \([N]\). Group the rational sum by left \(N(E^+)\)-cosets
and unfold to integrals of
\((\phi\otimes e)(x^{-1}n^{-1}\gamma y)\) over \(N(\mathbb A_{E^+})\).

Absolute convergence of the sum of their absolute integrals
follows from the compact-set bound for \(K\) with absolute values.
Writing \(\gamma\) for a representative of a left
\(N(E^+)\)-coset, the absolute convergence just proved allows
the product integral
\[
 \int_{N(\mathbb A_{E^+})}(\phi\otimes e)(x^{-1}n^{-1}\gamma y)\,dn
   =\prod_v\int_{N_v}(\phi\otimes e)_v(x_v^{-1}n_v^{-1}\gamma_v y_v)\,dn_v .
\]
The real constant-term factor is zero, so each coset integral
vanishes. The local groups here are the radicals of the
stabilizers of the corresponding isotropic subspaces. If there
is no isotropic subspace over \(E\), there is no proper rational
parabolic. Vanishing for maximal parabolics implies it for all
proper parabolics by successive integration of their radicals.

\proofstep{Pointwise vanishing of constant terms.}

Cuspidal descent needs pointwise vanishing for the chosen smooth vectors.


This constant-term vanishing extends from convolution of
compactly supported test functions to the full range of \(P\).
For a fixed compact set of \(x\)'s in the adelic group,
averaging \(F(nx)\) over a compact set of representatives of
\([N]\) defines a bounded map from \(L^2(Y)\) into local
\(L^2\) in \(x\). Write \(Q,T\) for these two compact sets. Integrate the pointwise
Cauchy--Schwarz bound over \(T\). With the fixed Haar measures,
\[
 \int_T\left|\int_{[N]}F(nx)\,dn\right|^2 dx
 \ll_Q \int_Q\int_T |F(nx)|^2\,dx\,dn
 \ll_Q \#\bigl(\Gamma\cap QT(QT)^{-1}\bigr)\,\|F\|_{L^2(Y)}^2 ,
\]
where \(\Gamma\) is the rational group.

The last estimate unfolds
each translate of \(T\) and bounds its overlap by the displayed
finite cardinality. Thus the vanishing is preserved by \(L^2\)
limits. The vectors \(F=T(w)\) have smooth representatives at
the real places by the local compact smoothing used for \(w_v\),
and are fixed by \(K_g\) at the finite places. For these
representatives the constant terms vanish everywhere by
continuity. Right translates, and hence their differentiated
vectors, retain the constant-term identities used in the
unipotent descent.
\end{proof}


\subsection{Small rational group elements}

The real projector suppresses elements far from the maximal compact. For rational elements satisfying the finite support condition, closeness at every real place forces all their eigenvalues to be one.

Set \(\epsilon_N=10^{-6}r^{-8}f^{-2}\).
\begin{lemma}[Small supported elements are unipotent]\label{lem:small-unipotent}
Suppose a supported rational element has a polar decomposition
\(U_v\exp X_v\) with \(\|X_v\|_{\rm HS}\le\epsilon_N\) at every
real place. It is unipotent and

\[
 \|U_v\exp X_v-I\|_{\rm HS}\le3\sqrt{r\epsilon_N}
\]
 at each such
place. If \(q=0\), the only supported rational element is the identity.
\end{lemma}

\begin{proof}
Since \(\|X_v\|_{\rm op}\le\epsilon_N\), for
\(1\le j\le r\) the trace of the \(j\)-th power differs from
\(r\) by at most \(r(1+\exp(j\epsilon_N))\le3r\).
The original rational element preserves the adelic lattice,
so its power traces are algebraic integers. Its reduction is
unipotent at every prime over \(g_0\). The traces of its powers
are therefore \(r\) modulo every such prime, and hence
modulo \(g_0\mathcal O_E\) because \(g_0\) is unramified.

Write \(a_j=\operatorname{tr}(\gamma^j)-r\).
If \(a_j\ne0\), divisibility and the archimedean estimates would give
\[
    g_0^{[E:\mathbb Q]}\le |N_{E/\mathbb Q}a_j|\le
        (3r)^{[E:\mathbb Q]}\qquad(1\le j\le r).
\]
This contradicts \(g_0>4r\).
Write \(p_j=\operatorname{tr}(\gamma^j)=r\) for \(1\le j\le r\).
Newton's recursion for the elementary symmetric functions gives
\[
 e_0=1,\qquad j e_j=\sum_{i=1}^j(-1)^{i-1} e_{j-i}p_i
      \quad(1\le j\le r).
\]
These are precisely the initial value and recursion of the identity
matrix, so \(e_j=\binom rj\). Hence \(\det(tI-\gamma)=(t-1)^r\).

Also \(\|U\exp X_v-U\|_{\mathrm{op}}\le2\epsilon_N\).
The equality of the first trace gives
\(\operatorname{Re}\operatorname{tr}U\ge r-2r\epsilon_N\).
Since \(U\) is unitary,
\(\|U-I\|_{\mathrm{HS}}^2=2r-2\operatorname{Re}\operatorname{tr}U\).
Subadditivity and \(\|A\|_{\rm HS}\le\sqrt r\|A\|_{\rm op}\) give
\[
 \|U\exp X_v-I\|_{\rm HS}
 \le 2\sqrt r\,\epsilon_N + (4r\epsilon_N)^{1/2}
 \le 3\sqrt{r\epsilon_N},
\]
using \(\epsilon_N\le1/4\).
This is the small entry bound used to place a nonidentity
unipotent in the canonical short isotropic flag of the lattice.
In the case \(q=0\), all archimedean elements lie in compact
groups and the power-trace argument applies to every element
in the support. A unitary unipotent is the identity by the
spectral theorem. The estimates for nonidentity elements below
are for \(q\ge1\).
\end{proof}


\subsection{Canonical lattice sizes and lowering by near elements}

Unipotence does not by itself bound how many rational translates can be near a given point. The subspaces of unusually small covolume identify a flag that constrains all such translates.

For a Hermitian lattice \(\Lambda\) at an adelic point,
give each rational subspace and quotient its induced
Euclidean metric at the real places. Write \(f=[E^+:\mathbb Q]\), with \(E^+\) the maximal totally real subfield of \(E\). The size of a nonzero lattice of rank \(j\) is its normalized covolume to the power \(1/(2fj)\). Normalize the real covolume by that of
\(\mathcal O_E^j\) in the standard product metric.

\begin{auxlemma}[The filtration by normalized covolume]
\label{lem:metric-size-filtration}
For an \(\mathcal O_E\)-lattice with the induced subspace and quotient
metrics, there is a unique filtration by saturated sublattices
whose grades are semistable and have strictly increasing sizes.
Here a lattice of size \(t\) is semistable if every nonzero
sublattice has size at least \(t\). If \(F(t)\) contains the grades
of size at most \(t\), then
\[
 \operatorname{size}(Q)\le t < \operatorname{size}(I)
\]
for every nonzero torsion-free quotient \(Q\) of \(F(t)\) and
every nonzero sublattice \(I\) of \(\Lambda/F(t)\), with the two
assertions understood separately. Duality reverses this filtration
and inverts the grade sizes.
\end{auxlemma}
\begin{proof}
\proofstep{Covolumes and attainment.}

For a saturated sublattice \(L'\subset L\) and the quotient
\(L''=L/L'\) with its induced metrics, normalized sizes satisfy
\[
 \operatorname{size}(L)^{\operatorname{rk}L}
 =\operatorname{size}(L')^{\operatorname{rk}L'}
  \operatorname{size}(L'')^{\operatorname{rk}L''}.
\]
This follows by using the determinant norm of the subspace
and the orthogonal projection to the quotient at the
infinite places. The same identity for the integral
determinants holds at every finite place. Saturating a
submodule can only decrease its size.

Among saturated sublattices of fixed rank there are only
finitely many with bounded size. In a fixed determinant ideal
class their exterior line lattices have the form
\(\mathfrak a w\) in the exterior power of \(\Lambda\), for
a fixed ideal representative \(\mathfrak a\). A bound for size
gives a bound for the product of the archimedean norms of \(w\).
Put \(b_v=\log\|w_v\|\) and \(\bar b=f^{-1}\sum_v b_v\).

Choose a unit \(u\) with
\[
     \max_v |\log |u|_v+b_v-\bar b|\le C_E .
\]
Here each complex embedding is taken once. The vector
\((\bar b-b_v)_v\) has sum zero, and \(C_E\) can be chosen as the
radius of a fundamental domain for the logarithms of the units. Dirichlet's unit lattice is
cocompact in the hyperplane of logarithms of product one.
Then \(uw\) generates the same exterior line lattice, since
\(u\) is a unit, and lies in a bounded set of the fixed fractional
exterior lattice and there are only finitely many possibilities.
If \(w\) represents that exterior line, the original
saturated sublattice is recovered by
\[
 L=\{v\in\Lambda:v\wedge w=0\}.
\]
Indeed over \(E\) this kernel is the subspace spanned by the
decomposable exterior vector, and saturation recovers its intersection
with \(\Lambda\).

\proofstep{The canonical increasing filtration.}

Let \(C(L)\) be the normalized covolume of a lattice \(L\),
with \(C(0)=1\) and \(C(L)=\operatorname{size}(L)^{2f\operatorname{rk}L}\).
For two saturated sublattices \(U,V\), put \(I=U\cap V\) and
\(W=(U+V)^{\rm sat}\). Project \(V\) orthogonally off \(U\).
The quotient \(V/I\) maps to the image lattice by a
volume-decreasing map, and saturation can only decrease the
volume further. The determinant formula in the two exact sequences
with quotients \(V/I\) and \(W/U\) therefore gives
\[
 C(I)C(W)\le C(U)C(V),\qquad
 \operatorname{rk}I+\operatorname{rk}W
                  =\operatorname{rk}U+\operatorname{rk}V .
\]
If \(U,V\) have the least size \(t\), put \(d_L=\operatorname{rk}L\).
Then
\[
 t^{2f(d_I+d_W)}\le C(I)C(W)\le C(U)C(V)
              =t^{2f(d_U+d_V)} .
\]
The two outer bounds coincide even when \(I=0\), with \(C(0)=1\).

In particular
\(W\) also has size \(t\). The maximal-rank minimizer is therefore
unique. In the quotient by that minimizer, a sublattice of size at
most \(t\) would have preimage of size at most \(t\) and larger
rank, since the preimage's covolume is the product of the two
factor covolumes. Thus the next step has strictly larger size. This constructs the increasing
Harder--Narasimhan sizes. Each grade of size \(t\) is
semistable, meaning that all its nonzero sublattices have
size at least \(t\). For a grade \(G\) of size \(t\) and a nonzero torsion-free
quotient \(Q=G/K\), multiplicativity and minimality give
\[
 C(Q)=\frac{C(G)}{C(K)}\le
       \frac{t^{2f\operatorname{rk}G}}{t^{2f\operatorname{rk}K}}
       =t^{2f\operatorname{rk}Q}.
\]
This includes \(K=0\).

To prove uniqueness, intersect a nonzero sublattice \(L\) with a
candidate filtration whose grade sizes are \(t_1<\cdots<t_a\).
Let \(d_i\) be the ranks of the induced grades of \(L\).
Each embeds integrally in its ambient grade. Orthogonal projection
to the ambient quotient can only decrease volume, as can saturation.
Multiplicativity on \(L\) therefore gives
\[
 C(L)\ge\prod_{i=1}^a t_i^{2fd_i}\ge t_1^{2f\operatorname{rk}L}.
\]

At equality all \(d_i\) for \(i>1\) vanish. The first step is thus
the maximal-rank minimizer, and induction gives uniqueness on its
quotient. Dualizing reverses the order of the
grades and takes reciprocal sizes, by the determinant
formula for a lattice and its integral dual. The quotient
size property of a semistable grade makes the dual grade
semistable. In particular if \(F(t)\) is the filtration
step with sizes at most \(t\), every nonzero torsion-free
quotient of \(F(t)\) has size at most \(t\), by duality and the
least-size inequality. Every nonzero sublattice of
\(\Lambda/F(t)\) has size strictly greater than \(t\).
\end{proof}

\begin{lemma}[The short flag of a self-dual Hermitian lattice]
\label{lem:canonical-short-flag}
The lattice has a canonical filtration with strictly increasing
sizes and semistable grades. Write \(F(t)\) for its step of sizes
at most \(t\). If the lattice is self-dual then its step \(S\) of
sizes strictly below \(1/2\) satisfies \(S^\perp=F(2)\),
\(\operatorname{rank}S\le q\), and the grades of \(S^\perp/S\)
have sizes in \([1/2,2]\).
If \(a\) is an integral endomorphism, functoriality in the size filtration gives 
\[
 \max_{v\mid\infty}\|a_v\|_{\rm op}\le\delta\quad\Longrightarrow\quad aF(t)\subseteq F(\delta t).
\]
 In particular, if
\(\|\gamma-I\|_{\rm op}\le\delta<1/4\) at all real places,
an integral unitary \(\gamma\) lies in the unipotent radical of
the stabilizer of the short flag. For an empty short flag
this forces \(\gamma=I\).
\end{lemma}

\begin{proof}
Use Lemma~\ref{lem:metric-size-filtration} with the normalized
covolume of the Hermitian lattice.
\proofstep{The perpendicular step and endpoint sizes.}

The Hermitian pairing makes the metric lattice self-dual.
Let \(S\) contain exactly the grades of size strictly less than \(1/2\).
Their dual grades have size strictly greater than \(2\). It follows that
\[
 S\subseteq S^\perp=F(2),\qquad \operatorname{rk}S\le q,\qquad
 \tfrac12\le\operatorname{size}(G)\le2
 \quad(G\text{ a grade of }S^\perp/S).
\] This describes the short flag and
fixes the treatment of the two endpoint sizes.

\begin{construction}
\begin{tikzpicture}[x=1cm,y=1cm,>=Stealth]
 \draw[->] (0,0)--(10.7,0) node[below left] {grade size};
 \draw[line width=.8pt] (3,-.08)--(3,.10) (7,-.08)--(7,.10);
 \node[below] at (3,-.06) {$1/2$};
 \node[below] at (7,-.06) {$2$};
 \node[above] at (1.4,.11) {$S$};
 \node[above] at (5,.11) {$S^\perp/S$};
 \node[above] at (8.9,.11) {$\Lambda/S^\perp$};
 \draw[<->] (1.45,.58) to[bend left=13]
   node[midway,above] {duality \(t\mapsto t^{-1}\)} (8.9,.58);
\end{tikzpicture}
\caption{The short filtration of Lemma~\ref{lem:canonical-short-flag}. Both boundary sizes belong to the middle quotient. The two outer groups of grades are dual.}\label{fig:short-reciprocal-grades}
\end{construction}

\proofstep{Lowering by a near element.}

The canonical filtration is intrinsic to the lattice. It constrains the difference of a near element from the identity, forcing it to lower the appropriate steps.


Suppose an integral endomorphism \(A\) of \(\Lambda\) has
operator norm at most \(\delta\) at every real place.
Let \(A_t:F(t)\to\Lambda/F(\delta t)\) be the induced map.
If its image has rank \(j>0\), pass to the saturated source
quotient \(Q_t\) and put \(I_t=(\operatorname{im}A_t)^{\rm sat}\)
in the target. The complex determinant norm is at most
\(\delta^j\) at each place. With the real squares and the
normalization of \(C\) this gives
\[
 \frac{C(I_t)}{C(Q_t)}\le\delta^{2fj},\qquad
 \operatorname{size}(I_t)\le
            \delta\operatorname{size}(Q_t)\le\delta t .
\]
Saturation accounts for the inequality on the left.
Every nonzero sublattice of the target has size strictly greater
than \(\delta t\), so \(A_t=0\).

Apply this with \(A=\gamma-I\) and the given \(0<\delta<1/4\).
The case \(\delta=0\) gives \(\gamma=I\) directly. It gives identity action
on each short grade, and the image of \(S^\perp=F(2)\)
under \(A\) lies in \(S\). Unitarity gives identity action
on the dual grades as well. Hence \(\gamma\) belongs to the
unipotent radical of the stabilizer of the short flag. If
the short flag is empty, this argument gives \(\gamma=I\).

\end{proof}

\subsection{Packing near elements on the flag blocks}

Lowering the short flag places the near elements in its unipotent radical. To count them one still needs the lattice spacing in each independent block of that radical.

\begin{lemma}[Near elements counted in the short-flag blocks]
\label{lem:near-flag-packing}
Use the grades of Lemma~\ref{lem:canonical-short-flag}.
For integral elements in the unipotent radical of the short flag
whose blocks of \(\gamma-I\) have Hilbert--Schmidt norms at
most \(\delta\) at every real place, the number is bounded by 
\[
 \prod_{v\mid\infty}\prod_{\text{free blocks at }v}(1+12\delta/\lambda)^d.
\]
 Here \(d\) is the block
dimension and
\(\lambda=\tau_i/\tau_j,\tau_i\tau_j,\tau_i,\tau_i^2\)
for the corresponding four block types in \(\mathcal L(\tau)\).
\end{lemma}

\begin{proof}
\proofstep{Induction by block distance.}

Complete the short flag by its perpendicular flag, using
the single middle block \(S^\perp/S\). Choose orthogonal
Euclidean complements at each real place, with dual blocks
paired by the unitary signature matrix. Write the resulting
blocks of \(\gamma-I\) in that order. They are strictly
upper triangular and each has Hilbert--Schmidt norm at
most \(\delta\). Count the possibilities in increasing
distance from the block diagonal. For elements agreeing on
all shorter distances, their difference maps each integral
filtration step into the step lower by the current distance.
Its blocks at that distance are therefore integral maps on
the graded lattices.

For two such elements write \(\gamma_i=1+A_i\) and
\(D=A_1-A_2\). Subtracting the two unitary identities gives
\[
      D+D^\dagger=-(D^\dagger A_1+A_2^\dagger D).
\]
The right side lowers by at least one more step, since
both \(A_i\) are strictly triangular. Thus at the current
distance the adjoint blocks in the difference are negatives
of one another's adjoints. It
suffices to count one block in each such pair and the
self-paired skew-Hermitian differences.

\proofstep{Block separation and volume packing.}

Once the successive free blocks are known, the remaining choices in each block form a separated set. This changes the problem from counting group elements to packing points in Euclidean balls.


For a nonzero difference in one block the same
kernel-to-saturated-image determinant argument bounds the
geometric mean of the operator norms below by the ratio
of the least target size to the largest source size.
For a nonzero block difference \(D\), write
\(l(D)=(\prod_v\|D_v\|_{\rm op})^{1/f}\). The resulting lower bounds are
\[
 l(D)\ge
 \begin{cases}
 \tau_i/\tau_j\ \text{or }\tau_i\tau_j,
                        &\text{the two cross blocks at }i<j,\\
 \tau_i/2,              &\text{a middle block and grade }i,\\
 \tau_i^2,              &\text{the self-paired block at }i .
 \end{cases}
\]
All the relevant Hilbert--Schmidt norms dominate the
operator norms. For an available ratio \(l\), put disjoint
product balls of radius \(l/2\) about the possible blocks
at the real places. An overlap would bound all difference
norms strictly below \(l\). In a self-paired block the
possibilities lie in an affine space with skew-Hermitian
direction, which is the space used for packing. Centering
at any one candidate bounds all radii from that center
by \(2\delta\).

If a finite subset has \(M\) candidates, translation and scaling
of product Lebesgue measure give
\[
 M(l/2)^d\operatorname{vol}(B_1)
       \le(2\delta+l/2)^d\operatorname{vol}(B_1),\qquad
 M\le(1+4\delta/l)^d ,
\]
where \(d\) includes all \(f\) places and \(B_1\) is the unit
product ball in the free block coordinates. This reasoning also
bounds an initially infinite candidate set, by applying the
volume inequality to any finite subset. Iterating through
the block distances determines the full matrices. The
dimensions of the free blocks are those in
\(\mathcal L(\tau)\) in \eqref{eq:count-flag-density-product}. Hence the number is at most
the product over those factors with \(\lambda^d\)
replaced by \((1+12\delta/\lambda)^d\), for
\(\lambda=\tau_i/\tau_j,\tau_i\tau_j,\tau_i,\tau_i^2\).
\end{proof}


\subsection{The measure of short flags}

Packing gives a bound at a fixed lattice. Its factors may grow when a grade of the flag becomes short, so the trace integral also needs the measure of lattices with those short grades.

Let \(\tau_1<\cdots<\tau_a<1/2\) be successive sizes of the
canonical short isotropic flag of Lemma~\ref{lem:canonical-short-flag}. Its layer ranks are positive
integers \(j_1,\ldots,j_a\), with \(J=\sum_i j_i\le q\).
\begin{lemma}[Measure of the short-flag region]\label{lem:short-flag-measure}
For \(q\ge1\) put \(U_0=rf\log r\). For large \(r\), let \(X_{\rm short}\) consist of the points in the quotient \(X\) with a nonempty isotropic flag of successive sizes less than \(1/2\). Its relative measure satisfies 
\[
 \frac{\operatorname{vol}(X_{\rm short})}{\operatorname{vol}(X)}\le e^{-U_0/4}.
\]

\end{lemma}

\begin{proof}
\proofstep{The modular character.}
The Levi group of its parabolic is the product of
\(\mathrm{GL}_{j_i}(E)\) and the unitary group of signature
\((p-J,q-J)\). In the logarithmic size variables the modular
character of the parabolic is the product of the determinants
on the unipotent blocks. Write this product as \(\mathcal L(\tau)\).
With \(f=[E^+:\Q]\) it is
\begin{align*}
 \mathcal L(\tau)
 &=\prod_{i<k}(\tau_i/\tau_k)^{2fj_i j_k}
               (\tau_i\tau_k)^{2fj_i j_k}\
 &\hspace{1em}\cdot\prod_i \tau_i^{2fj_i(r-2J)}
                                  (\tau_i^2)^{fj_i^2}.
\end{align*}
The two cross blocks between layers of complex dimensions
\(j_i,j_k\) each have real dimension \(2j_i j_k\) per place.
The blocks with the central quotient have real dimension
\(2j_i(r-2J)\). The skew-Hermitian block of layer \(i\) has
dimension \(j_i^2\) and scales by \(\tau_i^2\).

On the \(i\)-th layer of rank \(j_i\) use the positive determinant
coordinate \(a_i=\tau_i^{2fj_i}\). The quotient of its Haar
measure by the determinant-one gauge is \(da_i/a_i\).
The change from the determinant to the size on each layer is
\[
 \log a_i=2fj_i\log\tau_i,\qquad
 \prod_i\frac{da_i}{a_i}=\prod_i 2fj_i\frac{d\tau_i}{\tau_i}.
\]
The residual volumes below use this same decomposition.

Unfold the sum of short flags along this parabolic. The
normalized measure of any measurable set of sizes for the
given type is at most
\begin{equation}\label{eq:count-flag-density-product}
              \kappa\int\mathcal L(\tau)
                                   \prod_i \frac{2f j_i\,d\tau_i}{\tau_i}.
\end{equation}
The volume of the residual unitary quotient occurs in the
numerator of the unfolded integral. Dividing by the volume
of the original quotient replaces this factor by the inverse
of the successive unitary volume ratios. The ratios follow
from the convergent one-vector Siegel--Weil identity with
self-dual additive measures. In rank \(M\ge3\) average a
Gaussian sum over vectors of norm \(+1\) or \(-1\), using
the sign whose signature entry is positive. Witt's theorem
makes this a single orbit with stabilizer the group on
the perpendicular complement. Unfolding extracts the
volume of that rank-\((M-1)\) quotient.

On the Siegel--Weil
side the volume of the rank-\(M\) quotient multiplies
the local Gaussian densities and the Euler-product factor
\[
             d_E^{f/2}(d_E/2)^{-Mf}\mathcal L_M^{-1},
\]
where \(\mathcal L_M\) is the rank \(L\)-factor at \(M\).
The real Gaussian densities are positive for the selected
sign. The finite lattice normalization in the unfolding
contributes \(2^{Mf}\). Dividing this factor by the displayed
Euler-product factor gives the successive ratio in the common gauge,
\[
 \frac{2^{Mf}}{d_E^{f/2}(d_E/2)^{-Mf}\mathcal L_M^{-1}}
          =d_E^{f(M-1/2)}\mathcal L_M .
\]
Taking one step of each sign reduces the signature by \((1,1)\).
In the quotient of residual volume by original volume the resulting
Euler-product contribution after \(J\) pairs is
\[
 \prod_{a=0}^{J-1}(\mathcal L_{r-2a}\mathcal L_{r-2a-1})^{-1}
                 =\prod_{M=r-2J+1}^r \mathcal L_M^{-1}.
\] The general linear quotients give
\(\operatorname{Res}_{s=1}\zeta_E(s)\prod_{s=2}^{j_i}\zeta_E(s)\)
per layer. Haar gauges of the finite stabilizers and the
archimedean quotient occur in both the unfolding and the
volume ratio and cancel.

For clarity, in the resulting constant the discriminant
factor is \((d_E/2)^{-fd_N/2}\), where \(d_N\) is the sum
of the unipotent block dimensions at one place and
\(d_N\le2Jr\). The power of two is
\(2^{-f(rJ-J^2+J/2)}\). With \(S(m)=\prod_{i=1}^m
2\pi^i/(i-1)!\), the compact factor is
\[
 \left(\frac{S(p)S(q)}
              {S(p-J)S(q-J)\prod_i S(j_i)}\right)^f.
\]
The remaining factors are the general linear factors above
divided by the product of the rank \(L\)-factors at
\(r-2J+1,\ldots,r\), evaluated on the corresponding powers
of the quadratic character of \(E/E^+\).
The compact factor supplies the main saving in \(\kappa\).

For each \(p-J<i\le p\), one has \(i\ge r-2\sqrt r+1\),
and Stirling's estimate gives
\[
   \log\frac{2\pi^i}{(i-1)!}
            \le -r\log r+O(r+\sqrt r\log r)
\]
uniformly in that range. The other compact factors have
logarithms bounded above by \(O(fJ\sqrt r\log r)\), since
all their indices are at most \(\sqrt r\).

The discriminant factor costs at most \(O(fJr)\) in the
logarithm since \(d_E\ge3/2\), and the power of two is at most
one. A residue bound
\(\operatorname{Res}_{s=1}\zeta_E(s)\le
(16\log d_E)^{2f}\) follows from the unconditional degree
and discriminant bound for Dedekind zeta residues.
Since \(\log d_E\le e^r\), the residues also cost at most
\(O(fJr)\) in total. For the layer ranks \(j_i\), the Euler-product majorant gives
\[
 \sum_i\sum_{s=2}^{j_i}\log\zeta_E(s)
 \le 2f\,\#\{i\}\sum_{s\ge2}\log\zeta(s)\ll fJ .
\]
Indeed there are at most \(J\) layers, and \(\zeta(s)-1\ll2^{-s}\)
for integer \(s\ge2\).
For the reciprocal rank \(L\)-factors, the same Euler-product
bound in degree \(f\) gives \(O(fJ)\) in the logarithm.
The compact factors supply the leading negative term. Combining all
errors gives, uniformly for \(J\le q\le\sqrt r\) and large \(r\),
\[
                         \kappa\le \exp(-fJr\log r/2).
\]
\hfill\mbox{\cite{louboutin-row-count}}\label{cor:geometry-dimension-1}

\proofstep{Integration over the size variables.}

The modular character supplies powers of the relative sizes of the graded pieces. Integrating those powers determines whether the growth in the near-element count can be absorbed.

On the full simplex of sizes every factor of \(\mathcal L\)
is at most one. Keep just \(\prod_i\tau_i^{2fj_i(r-2J)}\).
When \(r-2J\ge1\), the integral in each remaining coordinate
is at most one, since
\[
 2f j_i\int_0^{1/2} t^{2fj_i(r-2J)-1}\,dt\le1 .
\]
Thus the relative measure of all nonempty short flags is
bounded by \(\sum_{J=1}^q2^{J-1}e^{-JU_0/2}\), where
\(U_0=rf\log r\). The sum is at most \(e^{-U_0/4}\) for
large \(r\). In particular the region \(\tau(x)<\tau_*<1/2\)
has this measure bound.
\end{proof}

\subsection{Absolute volume in the pointwise estimate}

The preceding estimate is relative to the total volume. The pointwise estimate in Section~\ref{sec:count-sobolev} instead uses the Haar density in Hilbert--Schmidt coordinates. To multiply the two bounds we compute the absolute quotient volume in that gauge.
\begin{lemma}[Absolute quotient volume in the local gauge]\label{lem:deep-volume-estimate}
At each real place let \(c_v\) be the reciprocal Haar density
at the origin in the Hilbert--Schmidt coordinates of
Section~\ref{sec:count-sobolev}. Use the quotient measure induced
by these Haar measures and the chosen finite measures. For \(X_0=rf d_Egb\), \(N=r^2 f\) in the stated
range (in particular \(r\ge8\), \(X_0\ge16\)),
\[
 \frac{(\prod_v c_v)V}{\operatorname{vol}(K_g)}
                 \le g^{2N}X_0^N.
\]
\end{lemma}

\begin{proof}
Use the unitary volume formula for a self-dual lattice over the
everywhere finitely unramified CM extension. If the finite stabilizer
and compact duals have volume one, the quotient volume in rank \(j\) is
\begin{equation}\label{eq:row-self-dual-mass}
  2 |D_{E^+}|^{j^2/2}
  \prod_{a=1}^j \left(\frac{(a-1)!}{(2\pi)^a}\right)^f L_a ,
\end{equation}
where \(L_a=L_{E^+}(a,\eta)\) for odd \(a\) and
\(L_a=\zeta_{E^+}(a)\) for even \(a\), with \(\eta\) the quadratic
character of \(E/E^+\). The Tamagawa number is \(2\). The local
densities at residue size \(Q\) are the orders of the corresponding
\(\mathrm{GL}_j\) or finite unitary groups divided by \(Q^{j^2}\).

If \(\eta_v=1\) at a split place and \(\eta_v=-1\) at an inert
one, both finite group order formulas read
\[
  Q^{-j^2}\#G(\mathbb F_Q)
        =\prod_{a=1}^j(1-\eta_v^{\,a}Q^{-a}).
\]
This is the inverse local Euler factor of \(\prod_{a=1}^j L_a\).
Self-duality at the unramified places supplies the smooth lattice
model, also in residue characteristic two. The real compact-dual
density is the same for the signatures compared here.
\cite{gan-yu-row-count,terakado-yu-row-count}
The volume of \(U(r)\) for the Hilbert--Schmidt metric is
\(\prod_{j=1}^r(2\pi)^j/(j-1)!\). To compute it, project to
the first column with fibre \(U(r-1)\).

At the identity a
horizontal tangent vector has first column \((ia,z)\), with
metric square \(a^2+2\|z\|^2\), while the unit sphere has
metric square \(a^2+\|z\|^2\). Thus the ratio of Hilbert--Schmidt volumes in successive ranks is
\[
 \frac{\operatorname{vol}U(r)}{\operatorname{vol}U(r-1)}
   =2^{r-1}\frac{2\pi^r}{(r-1)!}=\frac{(2\pi)^r}{(r-1)!}.
\]
Thus \(\operatorname{vol}U(r)=\prod_{j=1}^r(2\pi)^j/(j-1)!\). The passage from the compact dual
to \(U(p,q)\) multiplies the noncompact tangent directions by
\(i\), preserving their Hilbert--Schmidt measure. With the
gauge \(c_v\) defined in the Sobolev estimate, this gives
\[
  \frac{(\prod_v c_v)V}{\operatorname{vol}(K_f)}
       =2\,d_E^{r^2f/2}\prod_{j=1}^r L_j .
\]
Here the unramified CM extension has
\(|D_E|=|D_{E^+}|^2\), so \(|D_{E^+}|=d_E^f\).

The factors with \(j\ge2\) have product at most \(e^f\).
Indeed \(L_j\le\zeta(j)^f\) and
\(\sum_{j\ge2}\log\zeta(j)\le\sum_{n\ge2}1/(n(n-1))=1\).
For \(L_1=L_{E^+}(1,\eta)\), take the ratio of the
analytic class number formulas for \(E\) and \(E^+\).
The units of \(E^+\) give a sublattice in the unit lattice of
\(E\) of regulator \(2^{f-1}R_{E^+}\), since the complex
logarithmic coordinates are doubled. The regulator and the first Euler factor therefore satisfy
\[
 R_E\le2^{f-1}R_{E^+},\qquad
 L_1\le(2\pi)^f d_E^{-f/2}h_E .
\]

By Minkowski's bound every class of \(E\) has an ideal of
norm at most \(d_E^f=\sqrt{|D_E|}\). The ideal count up to
\(T\) is at most \(T(1+\log T)^{2f-1}\), by majorizing the
Euler coefficients by those of \(\zeta(s)^{2f}\) and
summing over the \(2f\) divisor factors. Hence
\[
 L_1\le (2\pi)^f d_E^{f/2}(1+f\log d_E)^{2f}.
\]
Put \(X_0=rf d_E g b\) and \(N=r^2 f\).

Since
\(1+f\log d_E\le f d_E\) and \(X_0\ge16\), the bound for
the quotient divided by \(\operatorname{vol}(K_f)\) is at most
\(X_0^{N/2+9f/2}\le X_0^N\) for \(r\ge8\).
Reduction at level \(g\) embeds \(K_f/K_g\) into the
endomorphisms of the rank-\(r\) lattice modulo \(g\).
There are \(g^{2fr^2}\) such endomorphisms, also for a projective
lattice by its local rank. It follows that
\[
   \frac{(\prod_v c_v)V}{\operatorname{vol}(K_g)}
                        \le g^{2N}X_0^N .
\]
This is the absolute bound used in the deep part of the
diagonal estimate.
\end{proof}

\subsection{Covering the short unipotent orbit}\label{sec:count-orbit}

In a deep part of the quotient, the lattice filtration has a large gap. The unipotent radical attached to this gap has a small set of representatives modulo its rational points, which allows cuspidality to control the value of a form.

\begin{proposition}[Covering in the chosen short flag]\label{prop:short-orbit-cover}
In Definition~\ref{hyp:cuspidal-dimension}, take \(q\ge1\) and suppose the smallest size \(\tau(x)\) in the canonical lattice filtration is less than \(1/2\). From the filtration in
Lemma~\ref{lem:canonical-short-flag} choose a step with consecutive
size ratio \(u\le(2\tau(x))^{1/q}\). Modulo its rational lattice
stabilizer and the finite level, the short unipotent orbit admits
representatives \(\exp A\exp Y\) in the real coordinates at \(x\),
where
\[
    \|A\|+\|Y\|\le \rho,\qquad
    \rho=\tfrac14 r^2\sqrt f\,g d_E(2\tau(x))^{1/q}.
\]
\end{proposition}

\begin{proof}
\proofstep{Choose a step by the product of ratios.}

Let \(\tau=t_1<t_2<\cdots<t_a<1/2\) be the short sizes
from Lemma~\ref{lem:canonical-short-flag} and put \(t_{a+1}=1/2\).
There are \(a\le q\) ratios since the sum of the short ranks
is at most \(q\). Then
\[
 \prod_{i=1}^a t_i/t_{i+1}=2\tau,\qquad
 \min_i t_i/t_{i+1}\le(2\tau)^{1/a}\le(2\tau)^{1/q}.
\]
Use a pair \(t,t'\) attaining one such ratio \(u\).
Let \(S\) be the saturated isotropic step ending at \(t\),
of rank \(j\). For the middle block \(Z=S^\perp/S\), the extreme sizes are
\[
 t_{\max}(S)=t,\qquad t_{\min}(Z)\ge t',\qquad t_{\min}(S^\vee)=1/t ,
\]
where the inequality for the middle block is needed only when it is nonzero.

\proofstep{Integral Witt coordinates and congruence lattices.}

Use an integral Witt splitting
\(\Lambda=S\oplus Z\oplus S^\vee\).
To obtain it, extend functionals on the saturated summand
\(S\) to \(\Lambda\) and use self-duality to identify
\(\Lambda/S^\perp\) with \(S^\vee\).
Choose a section by projectivity. Its Hermitian form can
be killed by subtracting an \(S\)-valued correction,
since the pairing between \(S\) and this section is perfect.
In detail, choose an integral scalar \(\theta\) with
\(\theta+\bar\theta=1\). It exists since the trace of the
unramified quadratic integral extension is surjective at
each prime, and so globally. Use a correction whose pairing
with the section is \(\theta\) times that section's form.
Indeed, if \(H_0\) is the Hermitian form of the section, the two
cross terms sum to
\[
                  \theta H_0+\bar\theta H_0=H_0 .
\]
Take the orthogonal complement of the resulting hyperbolic
pair to split off \(Z\).

The radical of the stabilizer of \(S\) has a two-step
Lie algebra with coordinates \(X:Z\longrightarrow S\)
and \(Y:S^\vee\longrightarrow S\), where \(Y\) is
skew-Hermitian. The other first-layer block is
\(-X^\dagger\). At the real places use the orthogonal Witt
coordinates. Two such splittings inducing the same maps on
the graded spaces differ by an element of the radical, as
their change of splitting is unitary, block triangular and
the identity on the graded spaces. Conjugation by it acts
trivially on the first-layer quotient and on the center.
Thus the congruence lattice supplies \(2g\operatorname{Hom}(Z,S)\)
in the first block.

For the first layer in the integral splitting, write
\[
 A_X=\begin{pmatrix}0&X&0\\0&0&-X^\dagger\\0&0&0\end{pmatrix},
 \qquad
 \exp(A_X)=
 \begin{pmatrix}1&X&-XX^\dagger/2\\
                 0&1&-X^\dagger\\0&0&1\end{pmatrix}.
\]
For \(X\in2g\operatorname{Hom}(Z,S)\) every off-diagonal
entry is integral and divisible by \(g\). In particular, writing
\(X=2gX_0\) makes the central entry
\[
 -XX^\dagger/2=-2g^2X_0X_0^\dagger
                   \in g\operatorname{Hom}(S^\vee,S).
\]
The center supplies \(g\) times the skew-Hermitian sublattice
of \(\operatorname{Hom}(S^\vee,S)\).

\proofstep{Dual minima of the two coordinate lattices.}

Here are the covering bounds for these lattices. In the
product Hilbert--Schmidt norm the Euclidean dual of
\(\operatorname{Hom}(Z,S)\) identifies with
\(2\mathfrak D_E^{-1}\operatorname{Hom}(S,Z)\), where
\(\mathfrak D_E\) is the absolute different. This follows
from the matrix trace pairing on projective modules and
the factor 2 relating field trace and the real product
pairing over the complex places. Multiplication of a
lattice by the codifferent multiplies every filtration
size by \(d_E^{-1}\), from its ideal norm and the same
exterior-volume normalization. The determinant comparison
for nonzero maps therefore bounds the geometric mean of
the local norms of a nonzero dual matrix below by
\(2d_E^{-1}t'/t\). For a nonzero vector \(B\) in this dual, the quadratic and
geometric means give
\[
 \Big(\sum_v\|B_v\|_{\rm HS}^2\Big)^{1/2}
 \ge\sqrt f\Big(\prod_v\|B_v\|_{\rm HS}\Big)^{1/f}
 \ge 2\sqrt f\,d_E^{-1}t'/t .
\] The first-layer lattice has
real dimension \(2fj(r-2j)\).

For the center, the skew-Hermitian sublattice is primitive
in the full matrix lattice. An integer functional on it
extends to the full lattice, so its Euclidean dual is the
orthogonal skew projection of the full dual. The projection \(B\mapsto (B-B^\dagger)/2\) places the dual of the
central lattice in \(\mathfrak D_E^{-1}\operatorname{Hom}(S,S^\vee)\).
The extreme sizes of \(S\) and its dual thus give
\[
                  \|B\|\ge \sqrt f\,d_E^{-1}/t^2
                         \qquad(B\ne0\text{ in the dual}).
\] The center has real dimension
\(fj^2\).

\proofstep{Cover each layer by transference.}

The dual minimum bounds give a covering radius in each additive layer. Combining these coverings must respect the order of the unipotent extension.


For a full Euclidean lattice of dimension \(d\), its
covering radius times the length of the shortest nonzero
dual vector is at most \(d/2\)
\cite[Theorem~2.2]{banaszczyk}. Apply this to the two
computed dual minima. Including the level \(2g\) in the
first layer and both its adjoint blocks, its radius is at most
\[
 \sqrt2\,(2g)\frac{j(r-2j)\sqrt f\,d_E}{2}\frac{t}{t'}
        \le\frac{r^2\sqrt f\,g d_E}{4\sqrt2}\frac{t}{t'} .
\]
In the center, including level \(g\), the radius is at most
\[
        \frac{g j^2\sqrt f\,d_E}{2}t^2
        \le\frac{r^2\sqrt f\,g d_E}{32}\frac{t}{t'} ,
\]
using \(t t'\le1/4\). First reduce the first block by an
integral congruence lift, then reduce the resulting center,
which already includes the commutator correction.
The two displayed radii can be added after this central correction.
For the gap \(t/t'\le(2\tau)^{1/q}\) this gives
\[
 \|A\|+\|Y\|\le r^2\sqrt f\,g d_E\frac{t}{t'}
       \left(\frac{1}{4\sqrt2}+\frac1{32}\right)
 \le \frac{r^2\sqrt f\,g d_E}{4}(2\tau)^{1/q}=:\rho .
\]

The radical is two-step with vector-group center and
vector-group quotient. Additive approximation at the
finite adeles in these two groups shows that each class
on its adelic orbit modulo the finite level has a real
representative. Apply the real covering just proved to
that representative. A rational left correction from
the lattice stabilizer at level \(g\) gives only a right
level translate at the point, so the value of each projected
function is unchanged. After conjugating into the real
coordinates at \(x\), the representative acts by a product
\(\exp A\exp Y\) with \(\|A\|+\|Y\|\le\rho\).
\end{proof}


\subsection{Bounds for the diagonal integral}

The diagonal is estimated separately where the smallest lattice size stays above the cutoff and where it falls below. In the latter region the covering of the short unipotent orbit is used before integration.

Write \(X_0=rf d_E g b\), \(N=r^2f\), \(\tau_*=X_0^{-11q}\).
The lattice size function \(\tau(x)\) is the threshold for its
first nonzero short layer. The portion \(\tau(x)\ge\tau_*\)
allows a uniform lattice point bound for the orbital expansion.
Write \(U_0=rf\log r\).
\begin{construction}
\begin{tabular}{@{}lll@{}}
\toprule
Region & Quantity controlled & Bound \\
\midrule
\(\tau(x)\ge\tau_*\) & integrated nonidentity terms &
       \((e^{-N}+e^{-U_0/4})D_*\)\\
\(\tau(x)<\tau_*\) & relative measure & \(e^{-U_0/4}\)\\
\(\tau(x)<\tau_*\) & integral for \(n\) orthonormal vectors & \(e^{-N}n\)\\
\bottomrule
\end{tabular}
\caption{Two estimates on the same quotient for \(q\ge1\) in Lemma~\ref{lem:diagonal-integral-bounds}. The final row uses an arbitrary finite family before taking the full trace.}\label{fig:two-trace-regions}
\end{construction}

\begin{lemma}[Bounds for the diagonal integral]\label{lem:diagonal-integral-bounds}
Suppose \(q\ge1\). If \(X\) is the quotient and \(X_{<}=\{x:\tau(x)<\tau_*\}\), the short-vector locus satisfies 
\[
 \frac{\operatorname{vol}(X_{<})}{\operatorname{vol}(X)}\le e^{-U_0/4}.
\]
 On its complement the integrated
absolute nonidentity terms are at most
\((e^{-N}+e^{-U_0/4})D_*\). On \(\tau(x)<\tau_*\) the integral of
the sum of squared values of any finite orthonormal family of
projected vectors is at most \(e^{-N}\) times the size of that family.
\end{lemma}

\begin{proof}
The measure estimate is Lemma~\ref{lem:short-flag-measure}.
\proofstep{Moderate region outside the polar threshold.}
  The point count for each archimedean shell has lattice factor
  at most \((3\sqrt r/\tau_*^2)^{2N}\).
  For two supported rational matrices their difference
  preserves the point lattice as an endomorphism. Thus two distinct supported matrices \(\gamma,\gamma'\) give a nonzero integral endomorphism \(T=\gamma-\gamma'\). The separation needed for packing is 
\[
 \left(\prod_v\|T_v\|_{\rm op}\right)^{1/f}\ge\tau_1(x)^2,\qquad \tau_1(x)=\min\{1/2,\tau(x)\},
\]
 with one representative from each complex-conjugate pair.
  In the size inequality for a nonzero lattice map, take
  the saturated quotient by the kernel and the saturated
  image. The determinant on the induced quotient has
  local norm at most the corresponding power of the
  operator norm. Write \(\Lambda_Q,\Lambda_I\) for the saturated quotient by
  the kernel and the saturated image. If their common rank
  is \(j\), the determinant inequality reads
  \[
   \operatorname{size}(\Lambda_I)
   \le \operatorname{size}(\Lambda_Q)
             \left(\prod_v\|A_v\|_{\rm op}\right)^{1/f}.
  \]

This follows by raising to \(2fj\), using the square of the
  complex determinant at each place, and allowing the index
  of the image in its saturation. The smallest size
  bounds \(\operatorname{size}(\Lambda_I)\) below and its
  reciprocal bounds \(\operatorname{size}(\Lambda_Q)\) above
  by self-duality. Their ratio is at least \(\tau_1(x)^2\).
  If the local containing balls have radii \(R_v\) of
  geometric mean \(R\), use product balls of radii
  \(\tau_1(x)^2 R_v/(2R)\) about the lattice points.
  The separation makes these product balls disjoint. Bounding their
total volume by that of the enlarged box gives
\[
    \#\{\text{points in the product ball}\}
       \le (1+2R/\tau_1(x)^2)^{2N}.
\]
  A noncompact distance at most \(u_v\) gives
  \(R_v\le\sqrt r\exp(u_v)\) by polar decomposition.
  Let \(\bar u=f^{-1}\sum_v u_v\). Since \(R\le\sqrt r e^{\bar u}\)
and \(\tau_1(x)\ge\tau_*\), the packing estimate gives
\[
 (1+2R/\tau_1(x)^2)^{2N}
 \le(1+2\sqrt r/\tau_*^2)^{2N}\exp(2r^2\textstyle\sum_v u_v).
\]

The large-gap matrix coefficient bound supplies a factor
  \(2f\,3^f\exp(-h\epsilon_N^2/16)\), which absorbs this
  lattice factor and \(e^N\). To see the shell sum put
  \(\epsilon=\epsilon_N\), \(c=\sqrt{2q}\) and \(B=2r^2\).
  If the noncompact distances at the places are bounded
  by \(u_v\), the lattice count is at most the lattice
  factor times \(\exp(B\sum_v u_v)\).
  For local noncompact distances \(s_v\), \eqref{eq:count-polar-decay}
bounds the coefficient relative to its identity value by
\[
 \prod_v\exp[-(h/8)\min(s_v^2,s_v/c)]
       =\exp[-(h/8)\textstyle\sum_v\min(s_v^2,s_v/c)] .
\]
This product permits a separate partition at each real place.
  Partition each distance into \([0,\epsilon]\) and
  \((2^{j-1}\epsilon,2^j\epsilon]\) for \(j\ge1\).
  For \(s_v\in[0,\epsilon]\) the counting factor costs at most
\(e^{B\epsilon}\le2\).

For \(u\le s_v\le2u\) the two competing
factors are bounded by
\[
 e^{Bs_v}\exp[-(h/8)\min(s_v^2,s_v/c)]
 \le \exp[2Bu-(h/8)\min(u^2,u/c)] .
\] For the
  bin with lower endpoint \(u=2^{j-1}\epsilon\), use
  \(\epsilon\le1/c\) and \(32B\le h\epsilon\) to get
  \[
   2Bu-\tfrac h8\min(u^2,u/c)\le-hu\epsilon/16,\qquad
   \sum_{j\ge1}e^{-h2^{j-1}\epsilon^2/16}
                         \le2e^{-h\epsilon^2/16}.
  \]
  Here \(h\epsilon^2/16\ge1\). Each coordinate costs at most
  \(3\) in total, and at least one must lie outside its first bin.
  Summing over that choice costs at most
  \(2f\,3^f e^{-h\epsilon^2/16}\). 

The lower bound supplied by the spectral gap and the logarithm of the lattice-counting factor are
  \[
    h\epsilon_N^2/16\ge50(C_g+4)r^3 f\log(bd_E),\qquad
    \log(\text{lattice factor})=2N(\log(3\sqrt r)+22q\log X_0).
  \]
  Here \(q\log X_0\le r+\sqrt r[\log r+
  (C_g+1)\log b+\log d_E]\). With \(b\ge2\) and
  \(d_E\ge3/2\), the spectral-gap bound absorbs this logarithm,
  \(N\) and \(\log(2f3^f)\), uniformly for large \(r\)
  at fixed \(C_g\).
\proofstep{Near unipotent terms in the moderate region.}
  Lemma~\ref{lem:small-unipotent} identifies the terms within the
  polar threshold as unipotent. Use Lemma~\ref{lem:near-flag-packing}
  to group the nonidentity terms by the ranks of the short flag. A type with \(a\) nonzero layers of
  total rank \(J\le q\) contributes at most
  \(2\kappa Y^a D_*\), where
  \(\kappa\le e^{-JU_0/2}\) and
  \(Y=2fq\log(1/\tau_*)\le e^{U_0/8}\).
  There are at most \(J\) layers, so its bound is
\[
 2\kappa Y^a D_*\le 2e^{(-J/2+a/8)U_0}D_*
              \le 2e^{-3JU_0/8}D_* .
\]
  The near count uses the same unipotent blocks as the
  measure. In a block of size parameter \(\lambda\) the count
  contributes \((1+12\epsilon/\lambda)^{d}\), where \(d\)
  is its real dimension and
  \(\epsilon=3\sqrt{r\epsilon_N}\). This is the lattice
  ball bound with radius \(\epsilon\) in the block gauge.
  Multiplying by the density changes each such factor to
  \((\lambda+12\epsilon)^d\).

Since \(\lambda\le1\) in each factor and \(\sum d\le3Jrf\), one has
\[
       \prod (\lambda+12\epsilon)^d
            \le\exp(12\epsilon\textstyle\sum d)
            \le e^{36Jrf\epsilon}\le2
\]
for sufficiently large \(r\).
  Integrating the remaining product of \(2fj_i\,d\tau_i/\tau_i\)
  from \(\tau_*\) to \(1/2\) in every coordinate costs at most
  \(Y^a\).
  A composition of \(J\) is determined by the cuts among its
\(J-1\) gaps. The total near contribution is bounded by
\[
  D_*\sum_{J=1}^q2^{J-1}\,2e^{-3JU_0/8}
      \le 2^{q+1}e^{-3U_0/8}D_* .
\] As \(q\le\sqrt r\), the near
  terms together contribute at most \(e^{-U_0/4}D_*\) for
  large \(r\).
\proofstep{Orthonormal families in the deep region.}

Deep in the quotient the short unipotent orbit bounds a whole orthonormal family by cuspidality.

  Use the fixed radius \(\rho\) from Proposition~\ref{prop:short-orbit-cover}
  at the point \(x\), and use the neighborhood of points
  \(x\exp(A)\exp(Y)\) of total coordinate norm at most \(\rho\).
  The deep threshold gives \(\rho\le\log2\), so the lattice
  size bound on this neighborhood is \(\tau(x)/2\).
  Write \(M_P\) for the supremum of the evaluation vector's
  norm on this neighborhood for unit vectors in the word
  space of order \(P\). Along each exponential segment the
  derivative inserts one infinitesimal action. The sum of
  the real-place derivative norms is at most \(\sqrt f B\)
  times the product norm of the direction, for
  \[
    B=\sqrt r\,b^2+\tfrac32 r^{3/2}+3(8qN+N).
  \]
  Each orbit representative is connected to \(x\) by two such
  segments of total length at most \(\rho\) within the
  neighborhood.

Let \(F_u(y)\) be the vector of evaluations of the family
  on a unit test vector \(u\) of order \(P\).
  Normalize the unipotent quotient measure to mass one and use
  cuspidality to write
  \[
    F_u(x)=\int_{[N]}\bigl(F_u(x)-F_u(nx)\bigr)\,dn .
  \]
  The real representatives from the orbit cover give the
  same values as \(nx\), by rational invariance and the finite
  level. At each representative the integral along the two
  exponential paths bounds the norm of the difference by
  \(\rho\sqrt f B M_{P+1}\). The same bound follows for
  \(F_u(x)\) after integration.

For any other point in the neighborhood the same segment
  estimate compares its value to the value at \(x\). The two comparisons give the recurrence on a single neighborhood
\[
             M_P\le2\rho\sqrt f B\,M_{P+1}\qquad(P<8qN).
\]
  The initial Sobolev estimate in Section~\ref{sec:count-sobolev}
  at order \(8qN\) gives the
  squared factor \((6B)^N e^f(17\sqrt r/\tau^2)^{2N}\)
  times the size of the family and the ratio of the archimedean gauge to the
  finite level volume. With \(K=2\rho\sqrt f B\), iterate the preceding inequality
on the same neighborhood to obtain
\[
 M_0^2\le K^{2j}M_j^2 \quad(0\le j\le8qN),\qquad
 M_0^2\le K^{16qN}M_{8qN}^2 .
\]
  In the same gauge the quotient volume is at most
  \(g^{2N}X_0^N\) times the reciprocal of that ratio.

  To check the resulting exponent, \(B\le X_0^4\) follows
  from \(q^2\le r\), \(b\ge2\), and \(X_0\ge16\).
  Then \(2\rho\sqrt f B\le X_0^6(2\tau)^{1/q}\).
  On taking the \(N\)-th root the initial bound and absolute
  volume give the factor
  \(6e\,17^2 B r\,\tau^{-4}g^2X_0\), with the gauge
  canceled. The descent gives the additional factor
  \((2\rho\sqrt f B)^{16q}\le
  X_0^{96q}(2\tau)^{16}\). Since \(B\le X_0^4\) and
  \(r g^2 X_0\le X_0^4\), their product is at most
  \[
     6e\cdot17^2\cdot 2^{16}X_0^{96q+8}\tau^{12}
       \le 6e\cdot17^2\cdot 2^{16}X_0^{-28}
       \le e^{-1}.
  \]
  In the first inequality use \(\tau\le X_0^{-11q}\).
  This gives the claimed diagonal integral.
\end{proof}

\subsection{Finite rank and the integral of the kernel}

Up to this point the estimates apply to every finite orthonormal family. Their uniformity now allows one to deduce finite rank and identify the integral of the diagonal with the dimension.

\begin{proof}[Proof of Proposition~\ref{prop:projector-dimension}]
\proofstep{Finite rank before taking a trace.}
The volume \(V\) is finite by Lemma~\ref{lem:deep-volume-estimate}. For
\(q\ge1\), apply the initial Sobolev bound at order zero on
\(\tau(x)\ge\tau_*\). For \(q=0\), the flag below size \(1/2\)
is empty by isotropy, so the same bound is uniform on the whole
quotient. In either case the bound gives one finite constant
\(C_2\) for the squared pointwise value of the unit lowest vector
under any isometric intertwiner on that region.

For an orthonormal family \(F_i=T_i(w)\), Schur's lemma and
evaluation on the unit vector \(w\) give \(T_i^*T_j=\delta_{ij}I\).
Put \(S(x)=\sum_{i=1}^m |F_i(x)|^2\). If \(S(x)>0\), the operator
\[
 T=\frac{\sum_i\overline{F_i(x)}T_i}{\sqrt{S(x)}}
 \quad\text{satisfies}\quad T^*T=I,\qquad T(w)(x)=\sqrt{S(x)} .
\]
The same pointwise bound therefore gives \(S(x)\le C_2\), also
when \(S(x)=0\). Integrate and apply
Lemma~\ref{lem:diagonal-integral-bounds} in the deep region:
\[
 \int S=m\le
 \begin{cases}C_2V+e^{-N}m,&q\ge1,\\ C_2V,&q=0.\end{cases}
\]
Every finite orthonormal family has bounded cardinality, so the
projection has finite rank.

\proofstep{Integration of the finite-rank kernel.}

Finite rank justifies computing the trace in an orthonormal basis. Its integral must still be identified with the estimate already obtained from periodization.

Now take an orthonormal basis \(F_1,\ldots,F_D\).
The integral kernel of this orthogonal projection is
\[
 K(x,y)=\sum_{j=1}^D F_j(x)\overline{F_j(y)}.
\]

Both this kernel and the periodized kernel are continuous
and represent the same bounded operator, so they agree
pointwise, as can be checked by testing on compactly supported
functions in the two variables. The integral of their diagonal
over \(Y\) is \(D\). For \(q=0\) the diagonal has only its
identity orbital term, giving \(D=D_*\).

For \(q\ge1\), the deep bound is now \(e^{-N}D\).
The identity has coefficient \(e(1)\) by Lemma~\ref{lem:finite-projector-mass}, and contributes \(D_*\) times the relative measure
of the moderate portion. The other moderate terms and the
deep contribution therefore give
\[
       \abs{D-D_*}
       \le (2e^{-U_0/4}+e^{-N})D_*+e^{-N}D .
\]
Write \(A=e^{-U_0/4}\) and \(B_1=e^{-N}\). Since
\(D\le D_*+\abs{D-D_*}\), move the last error to the left to
obtain \(\abs{D-D_*}\le 2(A+B_1)D_* /(1-B_1)\).
For large \(r\) this is at most \(8e^{-U_0/4}D_*\) and then
at most \(e^{-U_0/16}D_*\). All the thresholds used here
depend only on the fixed constants in the counting hypotheses.
This proves \eqref{eq:dimension-count}.
\end{proof}
